\documentclass[11pt,a4paper]{article}
\usepackage[margin=22mm]{geometry}
\usepackage{fontspec}
\usepackage[english]{babel}
\usepackage{xcolor}
\usepackage{graphicx}
\usepackage{amsmath,amssymb}
\usepackage{microtype}
\usepackage[hidelinks,bookmarksopen=true]{hyperref}
\usepackage{bookmark}
\setmainfont{DejaVu Sans}
\setlength{\parindent}{0pt}
\setlength{\parskip}{0.65em}
\definecolor{DeligneBlue}{HTML}{1D4E89}
\definecolor{LocatorGray}{HTML}{5E6770}
\hypersetup{pdftitle={D018: Restrained Critical Apparatus},pdfauthor={P. Deligne and M. Rapoport}}
\begin{document}
\begin{titlepage}\color{DeligneBlue}\vspace*{30mm}
{\fontsize{24}{30}\selectfont\bfseries D018: Restrained Critical Apparatus\par}
\vspace{8mm}{\Large P. Deligne and M. Rapoport\par}
\vspace{8mm}{\large Page-aligned authority, copy-matter, witness, and fallback notes\par}
\vfill{\small D018 --- source-aligned canonical edition\par}
\end{titlepage}
\tableofcontents\clearpage
\hypertarget{D018-P0001}{}
\pdfbookmark[2]{Authority page 1}{D018-P0001-bookmark}
\noindent{\color{LocatorGray}\footnotesize Authority page 1; collected-paper page 494; original folio 143.}\par\smallskip
Authority page 1; collected-paper volume page 494; original paper folio 143 (implicit on the unnumbered title leaf).

Included in the scholarly work: the exact title; the joint byline “par P. Deligne(1) et M. Rapoport(2)”; note (1), attached to Deligne; and note (2), attached to Rapoport. The joint authorship is controlling. The inherited metadata row naming only M. Rapoport is incomplete and was rejected as an authorship source.

Copy-matter disposition: “International Summer School on Modular Functions” and “Antwerp 1972” are preserved here as title-leaf provenance but are not repeated in either scholarly reader. There is no DeRa running locator on this page.

Witness comparison after the French record was frozen: the two diagram-aid branches contain byte-identical comparator-derived crops of the title, byline, and notes. Their legible wording agrees. The note crop also contains the Antwerp footer, which explains the deliberate body/provenance divergence. The immutable D018 identity receipt agrees on title, joint byline, note calls, venue line, authority hash, and boundary; it certifies no page transcription.

English: the title and byline are translated; note (1) remains as printed in English; note (2) is translated without changing its first-person institutional judgment.

Asset P0001-A01 preserves the underlined title, underlined joint byline, and the placement of note calls from the controlling authority pixels.


\par\medskip
\hypertarget{D018-P0002}{}
\pdfbookmark[2]{Authority page 2}{D018-P0002-bookmark}
\noindent{\color{LocatorGray}\footnotesize Authority page 2; collected-paper page 495; original folio 144.}\par\smallskip
Authority page 2; collected-paper volume page 495; original paper folio 144. Printed “-144-” and running locator “DeRa-2” are excluded from body prose and retained by this page address.

This is the first page of the SOMMAIRE. Entry order, wording, hierarchy, and printed internal page numbers are preserved. The source reads “Sous-schéma de non lissité” and “proreprésentabilité”; neither was silently normalized. Mathematical labels are transcribed as \ensuremath{\mathfrak{M}}\textunderscore{}*, \ensuremath{\mathfrak{M}}\textunderscore{}\{Γ₀(p)\}, \ensuremath{\mathfrak{M}}\textunderscore{}\{Γ₀₀(p)\}, and \ensuremath{\mathfrak{M}}\textunderscore{}n. The raw crop controls the exact fraktur form, subscripts, and alignment.

Witness comparison after freeze: seven page-specific inherited crops occur across the two byte-identical diagram-aid lineages. Legible entries agree with the frozen record. The Chapter V crop is detector-truncated at its line edges and therefore cannot serve as a complete notation witness. All inherited crops come from the lower-raster comparator and remain ZERO\textunderscore{}ACCEPTED.

English: section titles are rendered in standard algebraic-geometric terminology; the source pagination and distinction among the stack labels are retained.

Asset P0002-A01 preserves the lower contents block containing the special stack notation and right-aligned pagination.


\par\medskip
\hypertarget{D018-P0003}{}
\pdfbookmark[2]{Authority page 3}{D018-P0003-bookmark}
\noindent{\color{LocatorGray}\footnotesize Authority page 3; collected-paper page 496; original folio 145.}\par\smallskip
Authority page 3; collected-paper volume page 496; original paper folio 145. Printed “-145-” and running locator “DeRa-3” are excluded from body prose and retained by this page address.

This is the second and final contents page. The critical source distinction in VI.2 is between fraktur \ensuremath{\mathfrak{M}}\textunderscore{}H\textasciicircum{}∘ and roman M\textunderscore{}H\textasciicircum{}∘; it is not collapsed. The entries \ensuremath{\mathfrak{M}}\textunderscore{}\{Γ₀(p)\}, ℤ[[q]], and \ensuremath{\mathfrak{M}}\textunderscore{}H are also retained. “Pointes” is translated as “cusps” in the standalone English reader.

Witness comparison after freeze: the two diagram-aid branches provide byte-identical comparator-derived overinclusive crops of this contents block. No legible textual or pagination contradiction was found; the witness remains partial and non-authoritative.

Asset P0003-A01 preserves the complete printed contents block, including font distinction, subscripts, superscripts, and alignment.


\par\medskip
\hypertarget{D018-P0004}{}
\pdfbookmark[2]{Authority page 4}{D018-P0004-bookmark}
\noindent{\color{LocatorGray}\footnotesize Authority page 4; collected-paper page 497; original folio 146.}\par\smallskip
Authority page 4; collected-paper volume page 497; original paper folio 146. Printed “-146-” and running locator “DeRa-4” are excluded from body prose.

The Introduction heading, paragraph 1, all three displayed formula clusters, quotation marks around “à l'infini” and “structure de niveau”, and the page ending are included. The source alternates E\textunderscore{}z in the prose with E(z)\textunderscore{}n and E(z) in the torsion sentence; this notational distinction is preserved rather than harmonized. The final displayed isomorphism is transcribed in full as E(z)\textunderscore{}n ≃ (ℤ/nℤ)\textasciicircum{}2 : (a+bz)/n ↦ (a,b).

Witness comparison after freeze: the comparator-derived crops agree where legible. Inherited crop D018\textunderscore{}p004\textunderscore{}diag04 truncates the displayed isomorphism after its left portion and therefore diverges by omission; the authority pixels and P0004-A02 control the complete formula.

English: “transformations homographiques” is rendered “homographic transformations”; the modular and lattice notation is unchanged.

Assets P0004-A01 and P0004-A02 preserve, respectively, the half-plane/action displays and the full torsion isomorphism.


\par\medskip
\hypertarget{D018-P0005}{}
\pdfbookmark[2]{Authority page 5}{D018-P0005-bookmark}
\noindent{\color{LocatorGray}\footnotesize Authority page 5; collected-paper page 498; original folio 147.}\par\smallskip
Authority page 5; collected-paper volume page 498; original paper folio 147. Printed “-147-” and running locator “DeRa-5” are excluded from body prose.

The continuation from page 4, the definition of a level-n structure, the alternating form e\textunderscore{}n, ζ(α), ζ\textunderscore{}n = exp(2πi/n), the possible inverse sign convention, and paragraph 2 are included. The source wording “de caractéristique p ne divisant par n” is preserved in the frozen French; the English conveys the mathematical sense “characteristic p not dividing n” without altering the source record. Source line-end hyphenations “alter-/née” and “paramé-/trées” are retained in the diplomatic layer and recomposed in English.

Witness comparison after freeze: twelve inherited crops across the byte-identical branches cover overlapping fragments. All legible wording and formulas agree. Several crops are only sentence tails or detector fragments and have no independent semantic completeness; none is accepted as authority.

English: “structure de niveau n” is consistently rendered “level-n structure”; “familles ... paramétrées” is rendered as algebraic families parametrized by the base.

Assets P0005-A01 and P0005-A02 preserve the pairing/level-structure block and the cyclotomic normalization/sign-convention block.


\par\medskip
\hypertarget{D018-P0006}{}
\pdfbookmark[2]{Authority page 6}{D018-P0006-bookmark}
\noindent{\color{LocatorGray}\footnotesize Authority page 6; collected-paper page 499; original folio 148.}\par\smallskip
Authority page 6; collected-paper volume page 499; original paper folio 148. Printed “-148-” and running locator “DeRa-6” are excluded from body prose.

Parts (a)-(c) are included through the exact terminal words “du groupe des points”; the sentence continues on authority page 7 and is intentionally not completed here. Roman M\textunderscore{}n\textasciicircum{}∘ denotes the representing schemes on this page; it is not replaced by the fraktur stack notation used elsewhere. The functor map F\textunderscore{}nm → F\textunderscore{}n, Galois kernel, scalar extension, level-3 and level-4 coordinate rings, Hesse cubic, Legendre form, zero sections, and order-3 basis are preserved.

Witness comparison after freeze: eighteen inherited page-6 crops occur across the byte-identical lineages. Formula content agrees where legible. Several crops (notably diag03, diag05, and diag07-diag09) are isolated detector fragments; they diverge by truncation and cannot establish the page transition or a complete formula. The controlling authority assets below preserve each complete cluster.

English: “revêtement galoisien non ramifié” is rendered “unramified Galois covering”; “section neutre” is rendered “zero section”. The English page also ends mid-sentence.

Assets P0006-A01, P0006-A02, and P0006-A03 preserve the functor/covering block, the complete level-3 block, and the complete level-4 block, respectively.


\par\medskip
\hypertarget{D018-P0007}{}
\pdfbookmark[2]{Authority page 7}{D018-P0007-bookmark}
\noindent{\color{LocatorGray}\footnotesize Authority page 7; collected-paper page 500; original folio 149.}\par\smallskip
Authority page 7; collected-paper volume page 500; original paper folio 149. Printed “-149-” and running locator “DeRa-7” are excluded from body prose.

The page completes the order-4 construction from page 6, gives the displayed r and s formulas, states the gluing description of M\textunderscore{}n\textasciicircum{}∘[1/n], and begins the degeneration discussion. The page begins and ends mid-sentence. The printed subscript in π\textunderscore{}o(E'\textunderscore{}o) is o-shaped; the diplomatic record preserves it while this apparatus notes the mathematical component-group reading π\textunderscore{}0. In the pairing display the two inverse symbols are plain α\textasciicircum{}−1, while the specialized map is α\textunderscore{}o.

Witness comparison after freeze: all page-7 comparator-derived crops in both byte-identical diagram-aid lineages were inspected. Legible fragments agree; detector truncation and partial coverage prevent any authority role.

Assets P0007-A01 and P0007-A02 preserve the complete displayed formula and degeneration clusters from controlling authority pixels.


\par\medskip
\hypertarget{D018-P0008}{}
\pdfbookmark[2]{Authority page 8}{D018-P0008-bookmark}
\noindent{\color{LocatorGray}\footnotesize Authority page 8; collected-paper page 501; original folio 150.}\par\smallskip
Authority page 8; collected-paper volume page 501; original paper folio 150. Printed “-150-” and running locator “DeRa-8” are excluded from body prose.

The top tuple is (E\textunderscore{}o,e\textunderscore{}n,α\textunderscore{}n), whereas the later specialized structure is α\textunderscore{}o. The compactified quotient is barred X/Γ̄. Authority pixels distinguish the special fiber Ē'\textunderscore{}o from the minimal model Ē' with no subscript; both language records were corrected accordingly. The underlined term “courbe elliptique généralisée” is recorded in provenance.

The earlier P0008-A01 crop was rejected because it omitted the lower 92-pixel range of the unlabeled projective-line cycle. Its replacement is an exact one-bit crop around the complete diagram with a clean 16-pixel moat and no prose; P0008-A02 and P0008-A03 retain the surrounding formulas and definition data.

Witness comparison after freeze: all page-8 comparator crops in both byte-identical lineages were inspected as ZERO\textunderscore{}ACCEPTED only. Their legible fragments do not override the controlling authority.


\par\medskip
\hypertarget{D018-P0009}{}
\pdfbookmark[2]{Authority page 9}{D018-P0009-bookmark}
\noindent{\color{LocatorGray}\footnotesize Authority page 9; collected-paper page 502; original folio 151.}\par\smallskip
Authority page 9; collected-paper volume page 502; original paper folio 151. Printed “-151-” and running locator “DeRa-9” are excluded from body prose.

The barred compactified curve and functor X/Γ̄ and F̄\textunderscore{}n remain distinct from the later plain F\textunderscore{}n. Datum b) is the morphism + : E\textasciicircum{}reg ×\textunderscore{}T E → E; its source E\textasciicircum{}reg ×\textunderscore{}T E is not necessarily proper over T. The English wording was corrected so the source scheme is not misdescribed as a factor. The underlined source word “effective” is recorded in provenance.

Witness comparison after freeze: all page-9 comparator crops in both byte-identical lineages were inspected. Legible text and formulas agree, but the crops remain partial ZERO\textunderscore{}ACCEPTED evidence.

Assets P0009-A01 and P0009-A02 preserve the functorial and representability clusters.


\par\medskip
\hypertarget{D018-P0010}{}
\pdfbookmark[2]{Authority page 10}{D018-P0010-bookmark}
\noindent{\color{LocatorGray}\footnotesize Authority page 10; collected-paper page 503; original folio 152.}\par\smallskip
Authority page 10; collected-paper volume page 503; original paper folio 152. Printed “-152-” and running locator “DeRa-10” are excluded from body prose.

Both opening quotient occurrences are barred X/Γ̄; the prior minus-shaped transcription was repaired. The circular-looking affine-normalization sentence is preserved rather than silently regularized. Fraktur \ensuremath{\mathfrak{p}} denotes a prime ideal and plain p the rational prime. Braced notation M\textunderscore{}\{n'\}, M\textunderscore{}\{n'\}\textasciicircum{}∘, and M\textunderscore{}\{n''\}\textasciicircum{}∘ makes the printed n-prime and n-double-prime indices unambiguous without changing content. The page ends mid-sentence.

Witness comparison after freeze: both page-10 comparator crops in both byte-identical lineages were inspected as partial ZERO\textunderscore{}ACCEPTED evidence.

Assets P0010-A01 and P0010-A02 preserve the normalization formulas and reduction theorem continuation.


\par\medskip
\hypertarget{D018-P0011}{}
\pdfbookmark[2]{Authority page 11}{D018-P0011-bookmark}
\noindent{\color{LocatorGray}\footnotesize Authority page 11; collected-paper page 504; original folio 153.}\par\smallskip
Authority page 11; collected-paper volume page 504; original paper folio 153. Printed “-153-” and running locator “DeRa-11” are excluded from body prose.

The opening continuation uses M\textunderscore{}\{n'\}\textasciicircum{}∘ and M\textunderscore{}\{n''\}\textasciicircum{}∘; braced indices prevent confusion with primes attached to M\textunderscore{}n. The page preserves α\textasciicircum{}−1, β(p), M\textunderscore{}\ensuremath{\mathfrak{p}}, ℙ¹(ℤ/p\textasciicircum{}k), and the component index A. Fraktur \ensuremath{\mathfrak{p}} and plain p retain their different roles; the Pjateckii-Shapiro paragraph uses plain p. The page ends mid-sentence.

Witness comparison after freeze: all page-11 comparator crops in both byte-identical lineages were inspected. Legible fragments agree but remain partial ZERO\textunderscore{}ACCEPTED evidence.

Assets P0011-A01 and P0011-A02 preserve the p-primary morphism, component formulae, and continuation.


\par\medskip
\hypertarget{D018-P0012}{}
\pdfbookmark[2]{Authority page 12}{D018-P0012-bookmark}
\noindent{\color{LocatorGray}\footnotesize Authority page 12; collected-paper page 505; original folio 154.}\par\smallskip
Authority page 12; collected-paper volume page 505; original paper folio 154. Printed “-154-” and running locator “DeRa-12” are excluded from body prose.

The page completes the multiplicative level-structure formula, closes the supersingular-point discussion, and begins section 5. The underlined phrase “à distance finie” is translated “at finite distance.” Plain X/Γ is retained in the Fourier statements, while fraktur \ensuremath{\mathfrak{p}} denotes the prime ideal. The final bullet and closing sentence are complete on this page.

Witness comparison after freeze: all page-12 comparator crops in both byte-identical lineages were inspected. Legible content agrees; partial detector crops remain ZERO\textunderscore{}ACCEPTED and non-authoritative.

Assets P0012-A01 and P0012-A02 preserve the level-structure conclusion and Fourier-expansion results.


\par\medskip
\hypertarget{D018-P0013}{}
\pdfbookmark[2]{Authority page 13}{D018-P0013-bookmark}
\noindent{\color{LocatorGray}\footnotesize Authority page 13; collected-paper page 506; original folio 155.}\par\smallskip
Authority page 13; collected-paper volume page 506; original paper folio 155. Printed “-155-” and running locator “DeRa-13” are excluded from body prose.

The page states the Tate-curve heuristic principles, the Serre–Tate deformation argument, and the quaternionic modular-curve comparison. The barred finite field \ensuremath{\mathbb{F}}̄\textunderscore{}p, D\textunderscore{}p(E) = ⋃ E\textunderscore{}\{p\textasciicircum{}n\}, Δ, and X/Δ are retained. The two capital-A openings are unaccented in the authority and were refrozen accordingly.

Witness comparison after the initial French freeze: all four page-13 detector crops in both byte-identical lineages were inspected. Their legible content agrees; they remain partial ZERO\textunderscore{}ACCEPTED evidence with no authority role.

Assets P0013-A01 and P0013-A02 preserve the controlling authority pixels for the two main blocks.


\par\medskip
\hypertarget{D018-P0014}{}
\pdfbookmark[2]{Authority page 14}{D018-P0014-bookmark}
\noindent{\color{LocatorGray}\footnotesize Authority page 14; collected-paper page 507; original folio 156.}\par\smallskip
Authority page 14; collected-paper volume page 507; original paper folio 156. Printed “-156-” and running locator “DeRa-14” are excluded from body prose.

The mathematically unexpected printed expression \ensuremath{\mathcal{O}} ⊗\textunderscore{}ℤ ℤ\textunderscore{}p ≃ GL(2,ℤ\textunderscore{}p) is reproduced literally. The source-sic expression D\textunderscore{}p A) lacks a visible opening parenthesis and is also preserved rather than silently repaired. The page distinguishes X̄, K̄, X\textunderscore{}s̄, k(s̄), and ℚ̄, and ends mid-sentence after “posons.”

Witness comparison after the French record was frozen: all four page-14 crops in both byte-identical lineages were inspected. The full-page comparator crop and formula fragments corroborate the unusual GL reading and visible source anomaly but remain ZERO\textunderscore{}ACCEPTED.

Assets P0014-A01 and P0014-A02 preserve the p-divisible-group and restriction-of-scalars passages.


\par\medskip
\hypertarget{D018-P0015}{}
\pdfbookmark[2]{Authority page 15}{D018-P0015-bookmark}
\noindent{\color{LocatorGray}\footnotesize Authority page 15; collected-paper page 508; original folio 157.}\par\smallskip
Authority page 15; collected-paper volume page 508; original paper folio 157. Printed “-157-” and running locator “DeRa-15” are excluded from body prose.

The page completes the conjugate-scheme construction begun on page 14, then records the cyclotomic moduli morphism and both complex-base quotient identities. The printed ∑ is retained; the prose specifies that the sum is disjoint. The page ends after item 9a.

Witness comparison after freeze: all five page-15 crops in both byte-identical lineages were inspected. Legible formulas and prose agree but remain partial ZERO\textunderscore{}ACCEPTED evidence.

Assets P0015-A01 and P0015-A02 preserve the complete formula blocks and automorphism discussion.


\par\medskip
\hypertarget{D018-P0016}{}
\pdfbookmark[2]{Authority page 16}{D018-P0016-bookmark}
\noindent{\color{LocatorGray}\footnotesize Authority page 16; collected-paper page 509; original folio 158.}\par\smallskip
Authority page 16; collected-paper volume page 509; original paper folio 158. Printed “-158-” and running locator “DeRa-16” are excluded from body prose.

The page continues item 9b, states that M\textunderscore{}\{n̄m\} ramifies over M\textunderscore{}n at infinity, gives the coarse-moduli consequence, and closes the introduction with section 10 and the acknowledgment. The subscript is n̄m, with the bar over n alone. The large blank lower half is genuine page topology and has not been filled.

Witness comparison after freeze: both page-16 crops in both byte-identical lineages were inspected. They agree where legible and remain ZERO\textunderscore{}ACCEPTED.

Assets P0016-A01 and P0016-A02 preserve all substantive text while leaving the intentional blank region unaltered.


\par\medskip
\hypertarget{D018-P0017}{}
\pdfbookmark[2]{Authority page 17}{D018-P0017-bookmark}
\noindent{\color{LocatorGray}\footnotesize Authority page 17; collected-paper page 510; original folio 159.}\par\smallskip
Authority page 17; collected-paper volume page 510; original paper folio 159. Printed “-159-” and running locator “DeRa-17” are excluded from body prose.

This is the underlined Notations leaf. It preserves φ : \ensuremath{\mathfrak{M}}\textunderscore{}1 → \ensuremath{\mathfrak{M}}\textunderscore{}2, the representability fiber product, G\textunderscore{}n, base-change conventions, localization, and ζ\textunderscore{}n = exp(2πi/n). The word “représentable” is typographically underlined in the authority.

Witness comparison after freeze: all four page-17 crops in both byte-identical lineages were inspected. They corroborate the surrounding text but are incomplete and remain ZERO\textunderscore{}ACCEPTED.

Assets P0017-A01 and P0017-A02 preserve the stack and base-change notation from controlling pixels.


\par\medskip
\hypertarget{D018-P0018}{}
\pdfbookmark[2]{Authority page 18}{D018-P0018-bookmark}
\noindent{\color{LocatorGray}\footnotesize Authority page 18; collected-paper page 511; original folio 160.}\par\smallskip
Authority page 18; collected-paper volume page 511; original paper folio 160. Printed “-160-” and running locator “DeRa-18” are excluded from body prose.

The page begins Chapter I and defines curve schemes, arithmetic genus, and Cohen–Macaulay curve schemes. The authority prints C/S in the first paragraph of 1.1 but C|S in the next; the latter typography is preserved as source-sic rather than normalized. The formulas χ(\ensuremath{\mathcal{O}}\textunderscore{}\{C\textunderscore{}s\}) and 1 − g\textunderscore{}s = χ(C\textunderscore{}s,\ensuremath{\mathcal{O}}\textunderscore{}\{C\textunderscore{}s\}) retain their distinct arguments.

Witness comparison after freeze: all eight page-18 crops in both byte-identical lineages were inspected. Their legible content agrees, including C|S; all remain partial ZERO\textunderscore{}ACCEPTED evidence.

Assets P0018-A01 and P0018-A02 preserve the chapter opening and both definition blocks.


\par\medskip
\hypertarget{D018-P0019}{}
\pdfbookmark[2]{Authority page 19}{D018-P0019-bookmark}
\noindent{\color{LocatorGray}\footnotesize Authority page 19; collected-paper page 512; original folio 161.}\par\smallskip
Authority page 19; collected-paper volume page 512; original paper folio 161. Printed “-161-” and running locator “DeRa-19” are excluded from body prose.

The page opens the underlined section on coherent-sheaf duality, states the Cohen–Macaulay and Noetherian hypotheses, and records equations (2.1.1), (2.2.1), and (2.2.2). The derived Hom bases, degree shift [−d], and page-ending continuation are preserved.

Eight inherited crop witnesses across two lineages were hash-replayed from the preserved salvage archive after the French freeze. They form four byte-identical pairs and remain ZERO\textunderscore{}ACCEPTED comparison evidence with no authority role.

Assets P0019-A01 and P0019-A02 preserve the formula typography and page topology from controlling pixels.


\par\medskip
\hypertarget{D018-P0020}{}
\pdfbookmark[2]{Authority page 20}{D018-P0020-bookmark}
\noindent{\color{LocatorGray}\footnotesize Authority page 20; collected-paper page 513; original folio 162.}\par\smallskip
Authority page 20; collected-paper volume page 513; original paper folio 162. Printed “-162-” and running locator “DeRa-20” are excluded from body prose.

Equation (2.2.3) distinguishes the \ensuremath{\mathcal{O}}\textunderscore{}X-dual and \ensuremath{\mathcal{O}}\textunderscore{}S-dual. The page records the Gorenstein criterion and then characterizes regular differentials on a normalization by a residue sum, including the ordinary-double-point specialization. After π : X̃ → X, the authority unexpectedly prints “différentielles méromorphes sur X” without a tilde; that source-sic X is preserved.

Eighteen inherited crop witnesses across two lineages were hash-replayed after source freeze. They form nine byte-identical pairs and remain partial ZERO\textunderscore{}ACCEPTED evidence.

Assets P0020-A01 and P0020-A02 preserve both dense formula blocks.


\par\medskip
\hypertarget{D018-P0021}{}
\pdfbookmark[2]{Authority page 21}{D018-P0021-bookmark}
\noindent{\color{LocatorGray}\footnotesize Authority page 21; collected-paper page 514; original folio 163.}\par\smallskip
Authority page 21; collected-paper volume page 514; original paper folio 163. Printed “-163-” and running locator “DeRa-21” are excluded from body prose.

The page opens the relative Picard-functor section, defines fppf sheafification and cohomological flatness, and gives the degree formulas (3.3.1) and (3.3.2). The authority literally prints “Soit C un courbe propre.” Equation (3.3.2) is materially overstruck; its surviving final argument reads ℒ\textunderscore{}1 although ℒ\textunderscore{}2 is mathematically expected. Both anomalies are preserved rather than silently repaired.

Ten inherited crop witnesses form five hash-verified byte-identical lineage pairs. They remain ZERO\textunderscore{}ACCEPTED.

Assets P0021-A01 and P0021-A02 preserve the Picard and degree blocks.


\par\medskip
\hypertarget{D018-P0022}{}
\pdfbookmark[2]{Authority page 22}{D018-P0022-bookmark}
\noindent{\color{LocatorGray}\footnotesize Authority page 22; collected-paper page 515; original folio 164.}\par\smallskip
Authority page 22; collected-paper volume page 515; original paper folio 164. Printed “-164-” and running locator “DeRa-22” are excluded from body prose.

The page defines degree by irreducible-component multiplicities, the fppf degree morphism and Pic\textasciicircum{}\{[o]\}, then begins the unoriented component graph Γ(C). The typographically unusual prefixed-prime vertex classes 'Γ\textasciicircum{}o and "Γ\textasciicircum{}o, and the malformed printed item label (3.5 3), are retained.

Sixteen inherited crop witnesses form eight hash-verified pairs and remain ZERO\textunderscore{}ACCEPTED.

Assets P0022-A01 and P0022-A02 preserve the degree and graph-definition blocks.


\par\medskip
\hypertarget{D018-P0023}{}
\pdfbookmark[2]{Authority page 23}{D018-P0023-bookmark}
\noindent{\color{LocatorGray}\footnotesize Authority page 23; collected-paper page 516; original folio 165.}\par\smallskip
Authority page 23; collected-paper volume page 516; original paper folio 165. Printed “-165-” and running locator “DeRa-23” are excluded from body prose.

The page contains two hand-drawn curve/component-graph examples and a hand-drawn 2n-cycle. These are not replaced by approximate editable redrawing: authority pixels are carried in P0023-A01 and P0023-A02. The prose defines rotations through H\textasciicircum{}1(Γ(C),ℤ) and begins 3.7.

Sixteen inherited crop witnesses form eight byte-identical lineage pairs after hash replay and remain ZERO\textunderscore{}ACCEPTED.

The ○ and ⊗ vertex legend and the source's singular wording “celle” are documented without silent normalization.


\par\medskip
\hypertarget{D018-P0024}{}
\pdfbookmark[2]{Authority page 24}{D018-P0024-bookmark}
\noindent{\color{LocatorGray}\footnotesize Authority page 24; collected-paper page 517; original folio 166.}\par\smallskip
Authority page 24; collected-paper volume page 517; original paper folio 166. Printed “-166-” and running locator “DeRa-24” are excluded from body prose.

The page decomposes Pic\textunderscore{}\{C/k\}, records the total-degree map, the abelian quotient and torus kernel H\textasciicircum{}1(Γ(C),ℤ) ⊗ \ensuremath{\mathbb{G}}\textunderscore{}m, and the radicial comparison. The authority literally prints “et surjectif” after the total-degree display; this source-sic conjunction is preserved in French. It then opens the underlined non-smoothness section and ends mid-definition after Ω\textasciicircum{}1\textunderscore{}\{X/S\}.

Twenty-four inherited crop witnesses form twelve hash-verified lineage pairs and remain ZERO\textunderscore{}ACCEPTED.

Assets P0024-A01 and P0024-A02 preserve all displayed morphisms and the page transition.


\par\medskip
\hypertarget{D018-P0025}{}
\pdfbookmark[2]{Authority page 25}{D018-P0025-bookmark}
\noindent{\color{LocatorGray}\footnotesize Authority page 25; collected-paper page 518; original folio 167.}\par\smallskip
Authority page 25; collected-paper volume page 518; original paper folio 167. Printed “-167-” and running locator “DeRa-25” are excluded from body prose.

The page completes the Jacobian/non-smoothness discussion and opens deformation theory. It literally prints “Le lien où f est lisse,” uses the isomorphism arrow in (5.1.1), and has no visible period after X\textasciicircum{}reg. These surface features are retained or explicitly controlled in translation.

Fourteen inherited crop members form seven hash-verified byte-identical lineage pairs. They were replayed only after the source freeze and remain ZERO\textunderscore{}ACCEPTED comparison evidence with no authority role.

Assets P0025-A01 and P0025-A02 preserve the page transition, category notation, and formula typography.


\par\medskip
\hypertarget{D018-P0026}{}
\pdfbookmark[2]{Authority page 26}{D018-P0026-bookmark}
\noindent{\color{LocatorGray}\footnotesize Authority page 26; collected-paper page 519; original folio 168.}\par\smallskip
Authority page 26; collected-paper volume page 519; original paper folio 168. Printed “-168-” and running locator “DeRa-26” are excluded from body prose.

The versal node system uses an ordinary tensor product in the scalar-extension display. The round alpha subscript is not reliably distinguishable as o or 0, so the diplomatic text marks the uncertainty and the authority asset remains the controlling fallback. Both Proposition 5.3 and Théorème 5.3 occur literally. The proposition also lacks punctuation before “Alors.”

Eighteen inherited crop members form nine exact lineage pairs and remain ZERO\textunderscore{}ACCEPTED.

Assets P0026-A01 and P0026-A02 preserve the dense underlined statements and nodal equations.


\par\medskip
\hypertarget{D018-P0027}{}
\pdfbookmark[2]{Authority page 27}{D018-P0027-bookmark}
\noindent{\color{LocatorGray}\footnotesize Authority page 27; collected-paper page 520; original folio 169.}\par\smallskip
Authority page 27; collected-paper volume page 520; original paper folio 169. Printed “-169-” and running locator “DeRa-27” are excluded from body prose.

The page completes the nodal henselization theorem and gives the elliptic division-group discussion. The authority prints E\textunderscore{}n ≃ ℤ/n × μ\textunderscore{}n in the ordinary case and α\textunderscore{}\{p\textasciicircum{}2\} for the supersingular extension; neither is silently modernized. The exact sequence and its Cartier-duality interpretation are retained.

Sixteen inherited crop members form eight exact lineage pairs and remain ZERO\textunderscore{}ACCEPTED.

Assets P0027-A01 and P0027-A02 preserve all group-scheme displays and the physical page topology.


\par\medskip
\hypertarget{D018-P0028}{}
\pdfbookmark[2]{Authority page 28}{D018-P0028-bookmark}
\noindent{\color{LocatorGray}\footnotesize Authority page 28; collected-paper page 521; original folio 170.}\par\smallskip
Authority page 28; collected-paper volume page 521; original paper folio 170. Printed “-170-” and running locator “DeRa-28” are excluded from body prose.

Section 7 repeatedly prints the surface spelling “Cohen-MacAulay.” The displayed dimension formula is labeled “(7.1 1)” without the second dot, although the following prose cites “(7.1.1).” The diplomatic source preserves both forms; English uses the conventional proper-name spelling without changing the mathematics.

Eight inherited crop members form four hash-verified lineage pairs and remain ZERO\textunderscore{}ACCEPTED.

Assets P0028-A01 and P0028-A02 preserve the extensive underlining, formula, citations, and split page boundary.


\par\medskip
\hypertarget{D018-P0029}{}
\pdfbookmark[2]{Authority page 29}{D018-P0029-bookmark}
\noindent{\color{LocatorGray}\footnotesize Authority page 29; collected-paper page 522; original folio 171.}\par\smallskip
Authority page 29; collected-paper volume page 522; original paper folio 171. Printed “-171-” and running locator “DeRa-29” are excluded from body prose.

The page completes the fiberwise-property criterion and opens coarse moduli. Definition 8.1(ii) visibly uses plain s inside its parenthetical despite barred s elsewhere. In (8.2.1), the first stack-local object has a tilde-like mark over \ensuremath{\mathcal{O}}, whereas the coarse-space object has superscript h.

Six inherited crop members form three exact lineage pairs and remain ZERO\textunderscore{}ACCEPTED.

Assets P0029-A01 and P0029-A02 preserve the source-sic “surjectif,” the blackletter stack symbol, and the local-ring notation.


\par\medskip
\hypertarget{D018-P0030}{}
\pdfbookmark[2]{Authority page 30}{D018-P0030-bookmark}
\noindent{\color{LocatorGray}\footnotesize Authority page 30; collected-paper page 523; original folio 172.}\par\smallskip
Authority page 30; collected-paper volume page 523; original paper folio 172. Printed “-172-” and running locator “DeRa-30” are excluded from body prose.

The first local object now appears with superscript h. The page records the effective automorphism quotient, an étale presentation, and base-change behavior. It literally uses tensor symbols for stack/scheme base change, says “changements de base plat,” and ends with “est radicielle” without a period.

Six inherited crop members form three exact lineage pairs and remain ZERO\textunderscore{}ACCEPTED.

Assets P0030-A01 and P0030-A02 preserve the quotient isomorphism and base-change displays.


\par\medskip
\hypertarget{D018-P0031}{}
\pdfbookmark[2]{Authority page 31}{D018-P0031-bookmark}
\noindent{\color{LocatorGray}\footnotesize Authority page 31; collected-paper page 524; original folio 173.}\par\smallskip
Authority page 31; collected-paper volume page 524; original paper folio 173. Printed “-173-” and running locator “DeRa-31” are excluded from body prose.

The page opens Chapter II and constructs the standard Néron n-gon by cyclic gluing. The authority drawings show one self-nodal rational component for n=1 and three extended rational components forming a triangular cycle for n=3. The heading literally reads “Lemme 1 2.” with the dot between 1 and 2 absent.

Six inherited crop members form three hash-verified byte-identical lineage pairs. They were replayed only after source freeze and remain ZERO\textunderscore{}ACCEPTED comparison evidence with no authority role.

Assets P0031-A01 and P0031-A02 preserve the construction, both drawings, and the dualizing-sheaf formula.


\par\medskip
\hypertarget{D018-P0032}{}
\pdfbookmark[2]{Authority page 32}{D018-P0032-bookmark}
\noindent{\color{LocatorGray}\footnotesize Authority page 32; collected-paper page 525; original folio 174.}\par\smallskip
Authority page 32; collected-paper volume page 525; original paper folio 174. The printed folio ends in a detached upright stroke rather than a clean final hyphen; “DeRa-32” is excluded from body prose.

The page uses ω\textunderscore{}\{C/S\} in a lemma over an algebraically closed field, prints ω\textunderscore{}C ≃ \ensuremath{\mathcal{O}} without a visible subscript on \ensuremath{\mathcal{O}}, and ends with the literal but mathematically false inequality a > 0 ⇒ b\textunderscore{}i ≥ 1. The smooth solution on the next page has a=1 and b\textunderscore{}1=0, so the diplomatic editions preserve the source while this apparatus records the conflict.

Ten inherited crop members form five exact lineage pairs and remain ZERO\textunderscore{}ACCEPTED.

Assets P0032-A01 and P0032-A02 preserve dx/x, the residue signs, the overprinted normalization sequence, and all component-node inequalities.


\par\medskip
\hypertarget{D018-P0033}{}
\pdfbookmark[2]{Authority page 33}{D018-P0033-bookmark}
\noindent{\color{LocatorGray}\footnotesize Authority page 33; collected-paper page 526; original folio 175.}\par\smallskip
Authority page 33; collected-paper volume page 526; original paper folio 175. Printed “-175-” and running locator “DeRa-33” are excluded from body prose.

Definition 1.4 uses a special notion of stable genus-one curve. Proposition 1.5 literally prints both period and colon and uses singular “tel que.” Its proof also repeats “réduit” across several openness conditions. Proposition 1.6 and its base-change meaning are preserved without modernization.

Four inherited crop members form two exact lineage pairs and remain ZERO\textunderscore{}ACCEPTED.

Assets P0033-A01 and P0033-A02 preserve the numerical solutions, underlined statements, and equation (1.6.1).


\par\medskip
\hypertarget{D018-P0034}{}
\pdfbookmark[2]{Authority page 34}{D018-P0034-bookmark}
\noindent{\color{LocatorGray}\footnotesize Authority page 34; collected-paper page 527; original folio 176.}\par\smallskip
Authority page 34; collected-paper volume page 527; original paper folio 176. Printed “-176-” and the faint/damaged “DeRa-34” are excluded from body prose.

The page prints “I 3.7.” without the dot after I. More importantly, Lemma 1.8 literally begins with C̃ as the standard n-gon even though the proof immediately distinguishes the singular curve C from its normalization C̃. The source typo or stray tilde remains diplomatic; the mathematical distinction is not collapsed.

Six inherited crop members form three exact lineage pairs and remain ZERO\textunderscore{}ACCEPTED.

Assets P0034-A01 and P0034-A02 preserve the Picard formula, labelled exact sequence, normalization, and kernel action.


\par\medskip
\hypertarget{D018-P0035}{}
\pdfbookmark[2]{Authority page 35}{D018-P0035-bookmark}
\noindent{\color{LocatorGray}\footnotesize Authority page 35; collected-paper page 528; original folio 177.}\par\smallskip
Authority page 35; collected-paper volume page 528; original paper folio 177. Printed “-177-” and running locator “DeRa-35” are excluded from body prose.

The page gives the dihedral generators, the standard n-gon group action, and Proposition 1.10. Its semidirect-product symbol is visibly an x with a small raised prime-like mark; the diplomatic text records it as ×' while the prose supplies the semidirect-product meaning.

Ten inherited crop members form five exact lineage pairs and remain ZERO\textunderscore{}ACCEPTED.

Assets P0035-A01 and P0035-A02 preserve the automorphism displays, action (1.9.1), and complete proposition topology.


\par\medskip
\hypertarget{D018-P0036}{}
\pdfbookmark[2]{Authority page 36}{D018-P0036-bookmark}
\noindent{\color{LocatorGray}\footnotesize Authority page 36; collected-paper page 529; original folio 178.}\par\smallskip
Authority page 36; collected-paper volume page 529; original paper folio 178. Printed “-178-” and running locator “DeRa-36” are excluded from body prose.

The commutative diagram has an exact top row and three vertical arrows, while its lower row stops at Aut(\ensuremath{\mathbb{G}}\textunderscore{}m). Definition 1.12 prints the malformed label “(1 12.1),” clause a refers to the damaged and wrong “(1.1 .1),” and clause c changes from barred C\textunderscore{}s̄ to unbarred C\textunderscore{}s and Γ(C\textunderscore{}s). These anomalies are preserved and not silently harmonized.

Six inherited crop members form three exact lineage pairs and remain ZERO\textunderscore{}ACCEPTED.

Assets P0036-A01 and P0036-A02 preserve the diagram, character calculation, and full definition.


\par\medskip
\hypertarget{D018-P0037}{}
\pdfbookmark[2]{Authority page 37}{D018-P0037-bookmark}
\noindent{\color{LocatorGray}\footnotesize Authority page 37; collected-paper page 530; original folio 179.}\par\smallskip
Authority page 37; collected-paper volume page 530; original paper folio 179. Printed “-179-” and running locator “DeRa-37” are excluded from body prose.

Proposition 1.13 changes from geometric C\textunderscore{}s̄ to plain Pic⁰\textunderscore{}\{C\textunderscore{}s\} and literally says translation acts “par translation,” where rotation is mathematically expected. Lemma 1.14 never names V before “se factorise par V,” uses masculine “tel,” and later prints the malformed closure phrase V ⊃ V̄ est fermé. These readings are preserved rather than repaired.

Two inherited crop members form one hash-verified byte-identical lineage pair and remain ZERO\textunderscore{}ACCEPTED comparison evidence with no authority role.

Assets P0037-A01 and P0037-A02 preserve the propositions, Picard notation, formal-completion formula, and page break.


\par\medskip
\hypertarget{D018-P0038}{}
\pdfbookmark[2]{Authority page 38}{D018-P0038-bookmark}
\noindent{\color{LocatorGray}\footnotesize Authority page 38; collected-paper page 531; original folio 180.}\par\smallskip
Authority page 38; collected-paper volume page 531; original paper folio 180. Printed “-180-” and running locator “DeRa-38” are excluded from body prose.

The page varies barred and plain fiber notation: π₀(C\textunderscore{}s\textasciicircum{}reg) is plain, C\textunderscore{}reg/S has reg visibly lowered, and Lemma 1.16 changes from C\textunderscore{}s̄ to singular points of C\textunderscore{}s. The reference is literally “I 5.3,” and the line-broken ceux-ci has a faint initial c. No normalization is allowed to alter these source features.

Two inherited crop members form one exact lineage pair and remain ZERO\textunderscore{}ACCEPTED.

Assets P0038-A01 and P0038-A02 preserve the component calculation, lemma, local equation, and complete topology.


\par\medskip
\hypertarget{D018-P0039}{}
\pdfbookmark[2]{Authority page 39}{D018-P0039-bookmark}
\noindent{\color{LocatorGray}\footnotesize Authority page 39; collected-paper page 532; original folio 181.}\par\smallskip
Authority page 39; collected-paper volume page 532; original paper folio 181. Printed “-181-” and running locator “DeRa-39” are excluded from body prose.

After “Soit x ∈ U,” the proof literally says “près de U,” not près de x. It later prints “(ici, = V)” with nothing before the equals sign. The blowup, Stein factorization, and normalization formulas remain exact; the anomalies are recorded without supplying missing text.

Two inherited crop members form one exact lineage pair and remain ZERO\textunderscore{}ACCEPTED.

Assets P0039-A01 and P0039-A02 preserve the blowup argument and both anomalous proof lines.


\par\medskip
\hypertarget{D018-P0040}{}
\pdfbookmark[2]{Authority page 40}{D018-P0040-bookmark}
\noindent{\color{LocatorGray}\footnotesize Authority page 40; collected-paper page 533; original folio 182.}\par\smallskip
Authority page 40; collected-paper volume page 533; original paper folio 182. Printed “-182-” and running locator “DeRa-40” are excluded from body prose.

Diagram (1.17.1) visibly has X\textasciicircum{}reg ×\textunderscore{}S Y at bottom left, not X\textasciicircum{}reg ×\textunderscore{}S X. The proof first introduces “F' le foncteur” and later calls F' a subfunctor of F; u\textasciicircum{}\{−1\}ut(S) is also printed compactly. The two unlabeled-arrow diagrams are retained as authority assets and no hidden correction is made.

Two inherited crop members form one exact lineage pair and remain ZERO\textunderscore{}ACCEPTED.

Assets P0040-A01 and P0040-A02 preserve both diagrams, the unique-lift statement, and complete page transition.


\par\medskip
\hypertarget{D018-P0041}{}
\pdfbookmark[2]{Authority page 41}{D018-P0041-bookmark}
\noindent{\color{LocatorGray}\footnotesize Authority page 41; collected-paper page 534; original folio 183.}\par\smallskip
Authority page 41; collected-paper volume page 534; original paper folio 183. Printed “-183-” and running locator “DeRa-41” are excluded from body prose.

The constructed product is printed Y\textasciicircum{}reg × Y without \textunderscore{}S. The parenthesis after “x+y = y+x (x,y ∈ Y\textasciicircum{}reg(S)” never closes, and the final cyclic factor is printed ℤ/((n,m)) with doubled parentheses. All three features remain literal.

Four inherited crop members form two exact lineage pairs and remain ZERO\textunderscore{}ACCEPTED.

Assets P0041-A01 and P0041-A02 preserve the transported group law, multiplication-by-n map, and kernel formula.


\par\medskip
\hypertarget{D018-P0042}{}
\pdfbookmark[2]{Authority page 42}{D018-P0042-bookmark}
\noindent{\color{LocatorGray}\footnotesize Authority page 42; collected-paper page 535; original folio 184.}\par\smallskip
Authority page 42; collected-paper volume page 535; original paper folio 184. Printed “-184-” and running locator “DeRa-42” are excluded from body prose.

The rank at the top is printed n.(n,m). Corollary 1.20 uses n|m. Condition (2.1.1) is printed correctly in its display, while the immediate reference reads “(2 1.1)” without the first dot. The page ends after “pour”; no text from page 43 is supplied.

Six inherited crop members form three exact lineage pairs and remain ZERO\textunderscore{}ACCEPTED.

Assets P0042-A01 and P0042-A02 preserve Lemma 1.19, Corollary 1.20, section 2, and the complete displayed condition.


\par\medskip
\hypertarget{D018-P0043}{}
\pdfbookmark[2]{Authority page 43}{D018-P0043-bookmark}
\noindent{\color{LocatorGray}\footnotesize Authority page 43; collected-paper page 536; original folio 185.}\par\smallskip
Authority page 43; collected-paper volume page 536; original paper folio 185. Printed '-185-' and running locator 'DeRa-43' are excluded from body prose.

The page continues the proof of condition (2.1.1), gives the definition of m(t), states and completes Proposition 2.2, and proceeds through Lemma 2.2.1 and the Nakayama argument. The cross-reference is printed '(2 1.1)'. The line-break hyphen in 'isomor-/phe' is retained, and the actual page boundary occurs only after 'ce qui est'. The opening p\textunderscore{}*\ensuremath{\mathcal{O}}\textunderscore{}C line and EGA III. 7.8.8. reference are retained.

Four inherited diagram-crop members (two byte-identical lineage pairs) are replayed as ZERO\textunderscore{}ACCEPTED comparison evidence only. Assets P0043-A01 and P0043-A02 preserve the continuation, definition, proposition hypotheses, and exact page boundary.


\par\medskip
\hypertarget{D018-P0044}{}
\pdfbookmark[2]{Authority page 44}{D018-P0044-bookmark}
\noindent{\color{LocatorGray}\footnotesize Authority page 44; collected-paper page 537; original folio 186.}\par\smallskip
Authority page 44; collected-paper volume page 537; original paper folio 186. Printed '-186-' and running locator 'DeRa-44' are excluded from body prose.

The page closes the Nakayama argument for Proposition 2.2, states the cohomology equations (2.2.2) and (2.2.3), states Lemma 2.3, and includes the displayed duality consequence following 'On obtient'. The source-sic grammatical form 'le support de K est finie' and the source-sic reference '(2 1.1)' are retained. Page 45 begins with the parenthetical explanation of the dual symbol.

Two inherited diagram-crop members (one byte-identical lineage pair) are replayed as ZERO\textunderscore{}ACCEPTED comparison evidence only. Assets P0044-A01 and P0044-A02 preserve the exact sequence, lemma statement, and page transition.


\par\medskip
\hypertarget{D018-P0045}{}
\pdfbookmark[2]{Authority page 45}{D018-P0045-bookmark}
\noindent{\color{LocatorGray}\footnotesize Authority page 45; collected-paper page 538; original folio 187.}\par\smallskip
Authority page 45; collected-paper volume page 538; original paper folio 187. Printed '-187-' and running locator 'DeRa-45' are excluded from body prose.

The page completes the duality argument, beginning with the explanation of the printed dual symbol \textasciicircum{}∨, prints the heading 'Corollaire 2.4. :' with its visible punctuation, gives the universal Picard argument, and introduces the action and phi morphisms through (2.5.2). The action uses Pic⁰\textunderscore{}\{C/S\}, while (2.5.2) visibly targets Pic\textunderscore{}\{C/S\}; both are retained. The printed reference 'cf. I 3.2' and the source notation are retained; no later equation is supplied.

Fourteen inherited diagram-crop members (seven byte-identical lineage pairs) are replayed as ZERO\textunderscore{}ACCEPTED comparison evidence only. Assets P0045-A01 and P0045-A02 preserve the corollary, universal argument, and action definition.


\par\medskip
\hypertarget{D018-P0046}{}
\pdfbookmark[2]{Authority page 46}{D018-P0046-bookmark}
\noindent{\color{LocatorGray}\footnotesize Authority page 46; collected-paper page 539; original folio 188.}\par\smallskip
Authority page 46; collected-paper volume page 539; original paper folio 188. Printed '-188-' and running locator 'DeRa-46' are excluded from body prose.

The page begins with (2.5.3), states Theorem 2.6 and its three parts, gives the fppf Picard calculation, and begins Proposition 2.6.1. The proposition is deliberately cut after 'courbe intègre,' at the physical page boundary. Formulae and the distinction between phi and the Picard action are retained. The calculation's source-sic positive n\textunderscore{}i exponent inside the first bracket is preserved and is not silently normalized.

Eight inherited diagram-crop members (four byte-identical lineage pairs) are replayed as ZERO\textunderscore{}ACCEPTED comparison evidence only. Assets P0046-A01 and P0046-A02 preserve the theorem, proof, and proposition opening.


\par\medskip
\hypertarget{D018-P0047}{}
\pdfbookmark[2]{Authority page 47}{D018-P0047-bookmark}
\noindent{\color{LocatorGray}\footnotesize Authority page 47; collected-paper page 540; original folio 189.}\par\smallskip
Authority page 47; collected-paper volume page 540; original paper folio 189. Printed '-189-' and running locator 'DeRa-47' are excluded from body prose.

The page completes the torsion-free and Euler-characteristic argument of Proposition 2.6.1, states Proposition 2.7(i), and prints (2.7.1). The exact-sequence map labels f and e and the tensor exponent on m(e) are preserved. The local-on-S qualifier remains attached to the displayed formula and the record ends at the physical boundary; no text from page 48 is imported.

Four inherited diagram-crop members (two byte-identical lineage pairs) are replayed as ZERO\textunderscore{}ACCEPTED comparison evidence only. Assets P0047-A01 and P0047-A02 preserve the vanishing argument, quotient ideal, proposition, and formula.


\par\medskip
\hypertarget{D018-P0048}{}
\pdfbookmark[2]{Authority page 48}{D018-P0048-bookmark}
\noindent{\color{LocatorGray}\footnotesize Authority page 48; collected-paper page 541; original folio 190.}\par\smallskip
Authority page 48; collected-paper volume page 541; original paper folio 190. Printed '-190-' and running locator 'DeRa-48' are excluded from body prose.

The page states parts (ii)--(iv) of Proposition 2.7, gives (2.7.2), and continues the proof through the uniqueness argument. The distinct primed operation +′, the source-sic Hom term, the printed minus sign in the final φ\textunderscore{}e term (source-sic, not silently corrected), and the same-page line-break hyphenation 'locale-/ment sur S' are retained. The page ends after the completed phrase.

Eight inherited diagram-crop members (four byte-identical lineage pairs) are replayed as ZERO\textunderscore{}ACCEPTED comparison evidence only. Assets P0048-A01 and P0048-A02 preserve the proposition, formula, proof, and exact page boundary.


\par\medskip
\hypertarget{D018-P0049}{}
\pdfbookmark[2]{Authority page 49}{D018-P0049-bookmark}
\noindent{\color{LocatorGray}\footnotesize Authority page 49; collected-paper page 542; original folio 191.}\par\smallskip
Authority page 49; collected-paper volume page 542; original paper folio 191. Printed '-191-' and running locator 'DeRa-49' are excluded from body prose.

The page continues Proposition 2.7(iii), gives the two opening divisor identities, exact sequences (2.7.3)--(2.7.6), the coherent-duality calculation (2.7.5), and the degree/Euler-characteristic argument. The distinct printed forms \ensuremath{\mathcal{O}}((-t)) and \ensuremath{\mathcal{O}}(-(t)), every minus sign, and the phrase 'on a que' are retained. The physical page ends after 'On a donc bien'; page 50 completes that calculation.

Four inherited diagram-crop members (two byte-identical lineage pairs) are replayed as ZERO\textunderscore{}ACCEPTED comparison evidence only. Assets P0049-A01 and P0049-A02 are lossless authority crops preserving the dense formulae and exact page boundary; no reconstructed graphical content is introduced.


\par\medskip
\hypertarget{D018-P0050}{}
\pdfbookmark[2]{Authority page 50}{D018-P0050-bookmark}
\noindent{\color{LocatorGray}\footnotesize Authority page 50; collected-paper page 543; original folio 192.}\par\smallskip
Authority page 50; collected-paper volume page 543; original paper folio 192. Printed '-192-' and running locator 'DeRa-50' are excluded from body prose.

The page completes Proposition 2.7, states Corollary 2.8 and its fpqc/fppf proof, then opens Section 3 and subsection 3.1. The source visibly prints 'commutatifsplat' without a space; it is retained in the source-language edition and translated readably with this disclosure. The page ends after 'Soient les conditions suivantes.'

Six inherited diagram-crop members (three byte-identical lineage pairs) are replayed as ZERO\textunderscore{}ACCEPTED comparison evidence only. Assets P0050-A01 and P0050-A02 preserve the corollary, proof, section transition, and page boundary.


\par\medskip
\hypertarget{D018-P0051}{}
\pdfbookmark[2]{Authority page 51}{D018-P0051-bookmark}
\noindent{\color{LocatorGray}\footnotesize Authority page 51; collected-paper page 544; original folio 193.}\par\smallskip
Authority page 51; collected-paper volume page 544; original paper folio 193. Printed '-193-' and running locator 'DeRa-51' are excluded from body prose.

The page gives conditions (3.1.1)--(3.1.4), Theorem 3.2, its proof plan, and the reduction to assertions (A) and (B). The authority visibly prints S\textunderscore{}\{s̄\} in (3.1.4), the cross-reference '(3.1 1)' once without the second dot, and the source-sic verb form 'vérifié'; none is silently corrected in the frozen French record.

Four inherited diagram-crop members (two byte-identical lineage pairs) are replayed as ZERO\textunderscore{}ACCEPTED comparison evidence only. Assets P0051-A01 and P0051-A02 preserve the conditions, theorem, and proof plan.


\par\medskip
\hypertarget{D018-P0052}{}
\pdfbookmark[2]{Authority page 52}{D018-P0052-bookmark}
\noindent{\color{LocatorGray}\footnotesize Authority page 52; collected-paper page 545; original folio 194.}\par\smallskip
Authority page 52; collected-paper volume page 545; original paper folio 194. Printed '-194-' and running locator 'DeRa-52' are excluded from body prose.

The page proves assertion (A), begins the proof of Theorem 3.2 for G=ℤ, states and proves Lemma 3.3, and constructs g₁. Source-sic forms retained include g\textunderscore{}\{s̄\} ∈ G(s), omitted punctuation before 'On peut appliquer', 'Soient ζ', omitted punctuation after '(n fixe)', and '(3 1.3)'. The physical page ends mid-sentence after 'soit une translation'; page 53 begins '(cf. 1.17)'.

Four inherited diagram-crop members (two byte-identical lineage pairs) are replayed as ZERO\textunderscore{}ACCEPTED comparison evidence only. Assets P0052-A01 and P0052-A02 preserve the rotation argument, Lemma 3.3, construction of g₁, and exact boundary.


\par\medskip
\hypertarget{D018-P0053}{}
\pdfbookmark[2]{Authority page 53}{D018-P0053-bookmark}
\noindent{\color{LocatorGray}\footnotesize Authority page 53; collected-paper page 546; original folio 195.}\par\smallskip
Authority page 53; collected-paper volume page 546; original paper folio 195. Printed '-195-' and running locator 'DeRa-53' are excluded from body prose.

The opening parenthesis completes the sentence from page 52. The page then finishes the proof of Theorem 3.2 and gives Remark 3.4 with u : Pic\textasciicircum{}[0]\textunderscore{}\{C/S\} → C\textasciicircum{}reg and u(\ensuremath{\mathcal{O}}((t)-(e)))=t. Chapter II ends cleanly here. The large blank lower portion is genuine authority topology and is not treated as missing text.

Two inherited diagram-crop members (one byte-identical lineage pair) are replayed as ZERO\textunderscore{}ACCEPTED comparison evidence only. Assets P0053-A01 and P0053-A02 preserve the proof completion, remark, and blank-page topology.


\par\medskip
\hypertarget{D018-P0054}{}
\pdfbookmark[2]{Authority page 54}{D018-P0054-bookmark}
\noindent{\color{LocatorGray}\footnotesize Authority page 54; collected-paper page 547; original folio 196.}\par\smallskip
Authority page 54; collected-paper volume page 547; original paper folio 196. Printed '-196-' and running locator 'DeRa-54' are excluded from body prose. The pixel reading '-196-' controls over an erroneous OCR reading '-190-'.

The page begins Part III, 'Théorèmes de représentabilité', defines the stack \ensuremath{\mathfrak{M}} and substack \ensuremath{\mathfrak{M}}\textunderscore{}*, and begins the notation principles in 0.3. The Fraktur glyphs, fiber-product formula, superscripts o/h/∞, and p|n condition are retained. The physical page ends mid-sentence after 'de la courbe'; page 55 begins 'universelle.'

Two inherited diagram-crop members (one byte-identical lineage pair) are replayed as ZERO\textunderscore{}ACCEPTED comparison evidence only. Assets P0054-A01 and P0054-A02 preserve the chapter transition, stack notation, and exact page boundary.


\par\medskip
\hypertarget{D018-P0055}{}
\pdfbookmark[2]{Authority page 55}{D018-P0055-bookmark}
\noindent{\color{LocatorGray}\footnotesize Authority page 55; collected-paper page 548; original folio 197.}\par\smallskip
Authority page 55; collected-paper volume page 548; original paper folio 197. Printed '-197-' and running locator 'DeRa-55' are excluded from body prose.

The opening 'universelle.' completes notation item 0.3(d), whose page-54 text ends 'de la courbe'. This page gives notation items (e)--(f), opens Section 1, defines the deformation functor, and states Theorem 1.2. The stack symbol is visibly the lowercase fraktur \ensuremath{\mathfrak{m}}, distinct from coarse-moduli M. The authority also prints the source-sic 'Soient Λ un anneau'; it is retained in French and translated grammatically. The page ends with the deduction from (II.1.10) and [25].

Assets P0055-A01 and P0055-A02 are lossless authority crops preserving the boundary word, stack notation, theorem statement, and proof paragraph; no reconstructed graphical content is introduced.


\par\medskip
\hypertarget{D018-P0056}{}
\pdfbookmark[2]{Authority page 56}{D018-P0056-bookmark}
\noindent{\color{LocatorGray}\footnotesize Authority page 56; collected-paper page 549; original folio 198.}\par\smallskip
Authority page 56; collected-paper volume page 549; original paper folio 198. Printed '-198-' and running locator 'DeRa-56' are excluded from body prose.

The page treats the irreducible case, gives (1.3.1)--(1.3.2), and applies duality and normalization. In the duality display the Ext exponent is the variable i, paired with H exponent 1−i; this pixel reading controls over OCR. The source-sic phrase 'un algèbre de séries formelles' is retained. The later differential notation suppresses the superscript 1 exactly as printed. The page ends after 'on utilise que'; page 57 begins with the required cohomology display.

Assets P0056-A01 and P0056-A02 preserve the Ext formulas, duality, normalization morphism, and exact page boundary.


\par\medskip
\hypertarget{D018-P0057}{}
\pdfbookmark[2]{Authority page 57}{D018-P0057-bookmark}
\noindent{\color{LocatorGray}\footnotesize Authority page 57; collected-paper page 550; original folio 199.}\par\smallskip
Authority page 57; collected-paper volume page 550; original paper folio 199. Printed '-199-' and running locator 'DeRa-57' are excluded from body prose.

The opening display completes the page-56 cohomology argument. It visibly prints C and Ω\textunderscore{}C rather than the preceding C̃\textunderscore{}o notation and also prints an extra closing parenthesis after deg\textunderscore{}\{ℙ¹\}(e); both are retained in French rather than silently emended. The source also prints '(II.1.15.)' and the heading '1.4 Le cas général.' without a dot after 1.4. The page proves 1.2(iv), introduces D'', proves Lemma 1.4.1, and ends after 'on obtient'; page 58 begins 'ainsi un morphisme'.

Assets P0057-A01 and P0057-A02 preserve the anomalous opening formula, proof, functor definition, lemma, quotient construction, and exact boundary.


\par\medskip
\hypertarget{D018-P0058}{}
\pdfbookmark[2]{Authority page 58}{D018-P0058-bookmark}
\noindent{\color{LocatorGray}\footnotesize Authority page 58; collected-paper page 551; original folio 200.}\par\smallskip
Authority page 58; collected-paper volume page 551; original paper folio 200. Printed '-200-' and running locator 'DeRa-58' are excluded from body prose.

The opening words complete the page-57 sentence. The page proves Lemma 1.4.3, completes Theorem 1.2, opens Section 2 as 'Construction de \ensuremath{\mathfrak{m}}\textunderscore{}*', and proves the fpqc-stack assertion for \ensuremath{\mathfrak{m}} in Lemma 2.1 by descent. The stack glyph is lowercase fraktur, and the visibly doubled punctuation 'Lemme 2.1. :' is retained. The base strata S'\textunderscore{}n and S\textunderscore{}n and the curve subscheme C'\textunderscore{}n follow the authority. The page ends with the identification of 1-morphisms X → \ensuremath{\mathcal{G}} and objects of \ensuremath{\mathcal{G}}(X), citing [8].

Assets P0058-A01 and P0058-A02 preserve the proof transition, descent notation, citation, and lower page topology.


\par\medskip
\hypertarget{D018-P0059}{}
\pdfbookmark[2]{Authority page 59}{D018-P0059-bookmark}
\noindent{\color{LocatorGray}\footnotesize Authority page 59; collected-paper page 552; original folio 201.}\par\smallskip
Authority page 59; collected-paper volume page 552; original paper folio 201. Printed '-201-' and running locator 'DeRa-59' are excluded from body prose.

The page states Artin's algebraicity criterion as Theorem 2.3 and lists conditions (0)--(4). The limit display visibly uses the direct limit lim→\textunderscore{}i on the left and the inverse limit lim←\textunderscore{}i inside \ensuremath{\mathcal{G}}; both directions are frozen. The condition-(3) commutative square has Spec(k'\textunderscore{}o) → Spec(k\textunderscore{}o) across the top, vertical arrows u and ξ\textunderscore{}o, and Spec(R) → \ensuremath{\mathcal{G}} along the bottom with label ξ. The source punctuation 'catégories.).' is retained in French. The page ends with the complete neighborhood assertion in (4).

Assets P0059-A01 and P0059-A02 preserve the theorem and conditions. Asset P0059-A03 is an isolated lossless authority fallback for the commutative square, whose geometry cannot be fully represented by linear prose alone.


\par\medskip
\hypertarget{D018-P0060}{}
\pdfbookmark[2]{Authority page 60}{D018-P0060-bookmark}
\noindent{\color{LocatorGray}\footnotesize Authority page 60; collected-paper page 553; original folio 202.}\par\smallskip
Authority page 60; collected-paper volume page 553; original paper folio 202. Printed '-202-' and running locator 'DeRa-60' are excluded from body prose.

Remarks 2.4 replaces condition (3) by (3'a)--(3'b) under perfectness assumptions, defines Def(ξ\textunderscore{}o), and gives the normality/Krull-dimension simplification. Pixel inspection controls the hatted notation Ĉ\textunderscore{}\{Λ(k\textunderscore{}o)\}, barred ū, Aut isomorphism, and final source citation '([2]. Thm. 3.9)'. The source phrase 'classes d'isomorphies' is retained and translated idiomatically. The page ends cleanly; page 61 starts a new paragraph about the standard Néron n-gon and then Theorem 2.5.

Assets P0060-A01 and P0060-A02 preserve the deformation-category notation, conditions (3'a)--(3'b), and complete page ending; no diagram fallback is required.


\par\medskip
\hypertarget{D018-P0061}{}
\pdfbookmark[2]{Authority page 61}{D018-P0061-bookmark}
\noindent{\color{LocatorGray}\footnotesize Authority page 61; collected-paper page 554; original folio 203.}\par\smallskip
Authority page 61; collected-paper volume page 554; original paper folio 203. Printed '-203-' and running locator 'DeRa-61' are excluded from body prose.

The standard Néron n-gon map visibly has domain Spec(ℤ[1/n]); the comparator's Spec(ℤ) is rejected. The codomain star is visibly high in the opening formula, whereas Theorem 2.5 and its proof visibly set a low star and the non-smoothness locus adds ∞. The source condition (2) has the cyclic map φ : ℤ/n ↪ C₁, not a rank-two level structure, and maps Isom into an open of C\textunderscore{}n\textasciicircum{}reg. Assertion (*) retains the source-sic singular 'à n côté' and 'muni'. In the final line the authority omits the comma between the two arguments of Isom((C₁/H₁,e)(C₂/H₂,e)); this is frozen in French and normalized only in English.

Assets P0061-A01--P0061-A03 preserve the map, theorem, cyclic-subgroup construction, assertion (*), quotient-Isom morphism, and exact page ending.


\par\medskip
\hypertarget{D018-P0062}{}
\pdfbookmark[2]{Authority page 62}{D018-P0062-bookmark}
\noindent{\color{LocatorGray}\footnotesize Authority page 62; collected-paper page 555; original folio 204.}\par\smallskip
Authority page 62; collected-paper volume page 555; original paper folio 204. Printed '-204-' and running locator 'DeRa-62' are excluded from body prose.

The page completes conditions (3')--(4), formal smoothness, and Theorem 2.5, then gives Remark 2.6. The visibly run-together 'assertionde' and citation punctuation '1.2.(iii)', '2.5.', and 'II 1.15.' are source-sic and retained. The stack glyphs carry no h superscript: the open substacks are \ensuremath{\mathfrak{m}}\textunderscore{}1 and \ensuremath{\mathfrak{m}}\textunderscore{}(n)[1/n], inside \ensuremath{\mathfrak{m}}\textunderscore{}* and \ensuremath{\mathfrak{m}}\textunderscore{}*[1/n], respectively. The page ends with the complete comparison to \ensuremath{\mathfrak{m}}\textunderscore{}1[1/n].

Assets P0062-A01 and P0062-A02 preserve the proof conclusion, source-sic spacing, non-separatedness remark, and exact open-substack notation.


\par\medskip
\hypertarget{D018-P0063}{}
\pdfbookmark[2]{Authority page 63}{D018-P0063-bookmark}
\noindent{\color{LocatorGray}\footnotesize Authority page 63; collected-paper page 556; original folio 205.}\par\smallskip
Authority page 63; collected-paper volume page 556; original paper folio 205. Printed '-205-' and running locator 'DeRa-63' are excluded from body prose.

Chapter IV opens with the regime p∣n; the comparator's non-divisibility sign is rejected. Proposition 1.2 defines the contraction C → C̄. Its existence proof first gives H⁰(C\textunderscore{}s̄,\ensuremath{\mathcal{O}}(−nH\textunderscore{}s̄)) = 0 and only then derives H¹(C\textunderscore{}s̄,\ensuremath{\mathcal{O}}(nH\textunderscore{}s̄)) = 0 by duality. The homogeneous algebra display is visibly indexed n>0 even though the preceding sentence separately says that the same freeness assertion holds for n=0. The page ends after 'composante irréductible qui'.

Assets P0063-A01--P0063-A03 preserve the chapter opening, proposition, both vanishings, direct-sum index, and incomplete physical-page boundary.


\par\medskip
\hypertarget{D018-P0064}{}
\pdfbookmark[2]{Authority page 64}{D018-P0064-bookmark}
\noindent{\color{LocatorGray}\footnotesize Authority page 64; collected-paper page 557; original folio 206.}\par\smallskip
Authority page 64; collected-paper volume page 557; original paper folio 206. Printed '-206-' and running locator 'DeRa-64' are excluded from body prose.

The opening words and H⁰ isomorphism complete the page-63 contraction argument; the bars on H̄ and Ḡ over the contracted fiber are retained. The authority calls C̄ the homogeneous spectrum and later writes C̄ = Sph(Ḡ); the latter is retained rather than normalized to Proj. The uniqueness proof contains a triangle C → C̄ over S. Proposition 1.3 contains a second, square diagram: its top-left object is u⁻¹(C̄\textasciicircum{}reg) ×\textunderscore{}S C, its lower-left object is C̄\textasciicircum{}reg ×\textunderscore{}S C̄, its horizontal maps are +, and its vertical maps are u. These incidence relations are present in both the text layer and lossless fallback crops.

Assets P0064-A01--P0064-A03 preserve the opening isomorphism, finite-generation argument, triangle, proposition, and full square; P0064-A04 isolates the action square.


\par\medskip
\hypertarget{D018-P0065}{}
\pdfbookmark[2]{Authority page 65}{D018-P0065-bookmark}
\noindent{\color{LocatorGray}\footnotesize Authority page 65; collected-paper page 558; original folio 207.}\par\smallskip
Authority page 65; collected-paper volume page 558; original paper folio 207. Printed '-207-' and running locator 'DeRa-65' are excluded from body prose.

The page completes Proposition 1.3, gives Example 1.4, and opens the stable-reduction setup and Proposition 1.6. Formula (1.4.1) visibly runs (C\textasciicircum{}reg)\textasciicircum{}o →̃ c(C)\textasciicircum{}reg; the comparator's reversed arrow is rejected. The printed phrase 'D'après la démonstration de 1.3.4' is retained. In both the definition and Proposition 1.6(ii), the source visibly omits 'est' before 'à réduction stable'; English supplies ordinary grammar. Proposition 1.6(iii) says that the n-torsion points are defined over k(η), not over a barred generic point. The page ends at 'et il'.

Assets P0065-A01--P0065-A03 preserve the transported action, formula (1.4.1), stable-reduction wording, field k(η), divisibility n∣m, and exact boundary.


\par\medskip
\hypertarget{D018-P0066}{}
\pdfbookmark[2]{Authority page 66}{D018-P0066-bookmark}
\noindent{\color{LocatorGray}\footnotesize Authority page 66; collected-paper page 559; original folio 208.}\par\smallskip
Authority page 66; collected-paper volume page 559; original paper folio 208. Printed '-208-' and running locator 'DeRa-66' are excluded from body prose.

The opening completes Proposition 1.6(iii) and visibly prints the source-sic 'lissc'. Part (ii) extends the neutral section, translations, and inversion, then constructs the largest open U. The authority also prints 'résulterons'; both source-sic forms are retained and normalized only in English. The local-trait display is + : S₁ ×\textunderscore{}S C̃ → C̃. In the final paragraph the n-torsion group has exact rank n², the special fiber is smooth or an m-gon with n∣m, and the page closes after the contraction obtained from 1.3.

Assets P0066-A01--P0066-A03 preserve the proposition boundary, extension argument, local-trait display, torsion rank, and complete page ending.


\par\medskip
\hypertarget{D018-P0067}{}
\pdfbookmark[2]{Authority page 67}{D018-P0067-bookmark}
\noindent{\color{LocatorGray}\footnotesize Authority page 67; collected-paper page 560; original folio 209.}\par\smallskip
Authority page 67; collected-paper volume page 560; original paper folio 209. Printed '-209-' and running locator 'DeRa-67' are excluded from body prose.

The page completes Proposition 1.6(iv). The formal local equation is xy=t\textasciicircum{}k, and for k>1 the printed diagram replaces the node u by k−1 projective lines. The resolved curve C̃ has no exceptional curve of the first kind in C̃\textunderscore{}s and is the minimal model of C\textunderscore{}η; if it is not smooth, its special fiber is an nk-sided Néron polygon. Section 2 then defines \ensuremath{\mathfrak{m}}\textunderscore{}(n), with no comparator-added h superscript, and the section f\textunderscore{}n is defined by curve 1.9 over ℤ[1/n]. Proposition 2.2 prints \ensuremath{\mathfrak{m}}\textunderscore{}(n)\textasciicircum{}∞ = \ensuremath{\mathfrak{m}}\textunderscore{}(n) ∩ \ensuremath{\mathfrak{m}}\textunderscore{}*\textasciicircum{}∞ and the density of \ensuremath{\mathfrak{m}}\textunderscore{}(n)\textasciicircum{}o.

Assets P0067-A01--P0067-A03 preserve the proof continuation, minimal-model conclusion, section opening, and proposition; P0067-A04 isolates the resolution diagram.


\par\medskip
\hypertarget{D018-P0068}{}
\pdfbookmark[2]{Authority page 68}{D018-P0068-bookmark}
\noindent{\color{LocatorGray}\footnotesize Authority page 68; collected-paper page 561; original folio 210.}\par\smallskip
Authority page 68; collected-paper volume page 561; original paper folio 210. Printed '-210-' and running locator 'DeRa-68' are excluded from body prose.

The opening states that C\textunderscore{}n\textasciicircum{}reg is finite étale of exact rank n² and defines a level-n structure by α : C\textunderscore{}n →̃ (ℤ/n)². Definition 2.4 gives the stack \ensuremath{\mathfrak{m}}\textunderscore{}n[1/n] and its finite étale principal-homogeneous map to \ensuremath{\mathfrak{m}}\textunderscore{}(n). Its infinity locus is \ensuremath{\mathfrak{m}}\textunderscore{}n\textasciicircum{}∞[1/n]. Theorem 2.7 uses the printed citation ([26], app. à Exp. 17), not the altered comparator citation, and ends with faithful action of Aut(C) on C\textunderscore{}n.

Assets P0068-A01--P0068-A03 preserve the torsion rank, level map, stack definition, infinity locus, algebraic-space criterion, and page ending.


\par\medskip
\hypertarget{D018-P0069}{}
\pdfbookmark[2]{Authority page 69}{D018-P0069-bookmark}
\noindent{\color{LocatorGray}\footnotesize Authority page 69; collected-paper page 562; original folio 211.}\par\smallskip
Authority page 69; collected-paper volume page 562; original paper folio 211. Printed '-211-' and running locator 'DeRa-69' are excluded from body prose.

Variant 2.8 gives the fpqc sheaf F\textunderscore{}n and Artin's criterion. Corollary 2.9 identifies \ensuremath{\mathfrak{m}}\textunderscore{}n[1/n] as a smooth projective scheme, with projectivity supplied by the relatively ample divisor \ensuremath{\mathfrak{m}}\textunderscore{}n\textasciicircum{}∞. Section 3 defines F\textunderscore{}H from the left action of H. The source itself warns that, because H acts on the left, H\textbackslash{}F(T) might be preferable to the printed F(T)/H; both the quotient and the source-sic spacing gauche,il are retained. The page ends after F\textunderscore{}H = Hom(C\textunderscore{}n,(ℤ/n)²)/H.

Assets P0069-A01--P0069-A03 preserve the representability route, ample divisor, quotient warning, and exact terminal Hom formula.


\par\medskip
\hypertarget{D018-P0070}{}
\pdfbookmark[2]{Authority page 70}{D018-P0070-bookmark}
\noindent{\color{LocatorGray}\footnotesize Authority page 70; collected-paper page 563; original folio 212.}\par\smallskip
Authority page 70; collected-paper volume page 563; original paper folio 212. Printed '-212-' and running locator 'DeRa-70' are excluded from body prose.

The page explains descent for level-H structures: local representatives α and β differ by a uniquely determined g∈H. Definition 3.2 uses the o-open stack \ensuremath{\mathfrak{m}}\textunderscore{}H\textasciicircum{}o[1/n], finite étale and locally representable over \ensuremath{\mathfrak{m}}\textunderscore{}1\textasciicircum{}o[1/n]. Definition 3.3 makes \ensuremath{\mathfrak{m}}\textunderscore{}H the normalization of \ensuremath{\mathfrak{m}}\textunderscore{}1 in that open stack, with no prime on \ensuremath{\mathfrak{m}}\textunderscore{}1. Theorem 3.4 distinguishes \ensuremath{\mathfrak{m}}\textunderscore{}H[1/n], \ensuremath{\mathfrak{m}}\textunderscore{}H\textasciicircum{}o[1/n], and the finite étale boundary \ensuremath{\mathfrak{m}}\textunderscore{}H\textasciicircum{}∞[1/n]. The source forms Abhyankhar and appliquable are frozen in French.

Assets P0070-A01--P0070-A03 preserve the local quotient description, exact stack superscripts, normalization, theorem, and proof.


\par\medskip
\hypertarget{D018-P0071}{}
\pdfbookmark[2]{Authority page 71}{D018-P0071-bookmark}
\noindent{\color{LocatorGray}\footnotesize Authority page 71; collected-paper page 564; original folio 213.}\par\smallskip
Authority page 71; collected-paper volume page 564; original paper folio 213. Printed '-213-' and running locator 'DeRa-71' are excluded from body prose.

For H=\{e\}, the authority sets \ensuremath{\mathfrak{m}}\textunderscore{}n=\ensuremath{\mathfrak{m}}\textunderscore{}H. Proposition 3.5 identifies the full stacks after inverting n, while its proof separately identifies \ensuremath{\mathfrak{m}}\textunderscore{}n\textasciicircum{}o[1/n] with \ensuremath{\mathfrak{m}}\textunderscore{}H\textasciicircum{}o[1/n]. Lemma 3.5.1 states the injectivity of Aut(C)→Aut(c(C)). In 3.6, H' is the inverse image of H in GL(2,ℤ/nm), and the displayed square prints x\textasciicircum{}m for the multiplication-by-m map on E\textunderscore{}nm and réduction for the constant-group map. The physical page ends with \ensuremath{\mathfrak{m}}\textunderscore{}H'\textasciicircum{}o[1/nm] ≃ \ensuremath{\mathfrak{m}}\textunderscore{}H\textasciicircum{}o[1/n][1/nm].

Assets P0071-A01--P0071-A03 preserve the proposition, contraction map, lemma, torsion square, and page-ending formula; P0071-A04 isolates the square.


\par\medskip
\hypertarget{D018-P0072}{}
\pdfbookmark[2]{Authority page 72}{D018-P0072-bookmark}
\noindent{\color{LocatorGray}\footnotesize Authority page 72; collected-paper page 565; original folio 214.}\par\smallskip
Authority page 72; collected-paper volume page 565; original paper folio 214. Printed '-214-' and running locator 'DeRa-72' are excluded from body prose.

The opening completes 3.6 by identifying \ensuremath{\mathfrak{m}}\textunderscore{}H'≃\ensuremath{\mathfrak{m}}\textunderscore{}H and then denotes the stack by the inverse-image open subgroup K⊂GL(2,ℤ̂). Section 3.7 uses coprime n and m and the reduction isomorphism GL(2,ℤ/nm)→̃GL(2,ℤ/n)×GL(2,ℤ/m). Construction 3.8 gives the fiber product over \ensuremath{\mathfrak{m}}\textunderscore{}1\textasciicircum{}o. Formula (3.8.1) is C\textunderscore{}nm→̃C\textunderscore{}n×C\textunderscore{}m. The barred symbols ᾱ and β̄ are structure classes, while α and β are local representatives; both distinctions and the reduction square are retained.

Assets P0072-A01--P0072-A03 preserve the normalization, group decomposition, stack isomorphism, torsion product, and final map; P0072-A04 isolates the commutative square.


\par\medskip
\hypertarget{D018-P0073}{}
\pdfbookmark[2]{Authority page 73}{D018-P0073-bookmark}
\noindent{\color{LocatorGray}\footnotesize Authority page 73; collected-paper page 566; original folio 215.}\par\smallskip
Authority page 73; collected-paper volume page 566; original paper folio 215. Printed '-215-' and running locator 'DeRa-73' are excluded from body prose.

Proposition 3.9 extends Construction 3.8 from the open loci to the compactified fiber product \ensuremath{\mathfrak{m}}\textunderscore{}K×\textunderscore{}\{\ensuremath{\mathfrak{m}}\textunderscore{}1\}\ensuremath{\mathfrak{m}}\textunderscore{}L. Both fiber-product bases and the intersection K∩L are frozen from the authority. The proof uses the two localizations asymmetrically and then invokes tame ramification and Serre's criterion. The printed Corollary 3.9.2 says 'et ≥ 3' without repeating n. Proposition 3.10 begins the systematic distinction between the stack \ensuremath{\mathfrak{m}}\textunderscore{}H and its coarse space M\textunderscore{}H; the last source sentence also visibly passes directly from \ensuremath{\mathfrak{m}}\textunderscore{}H→\ensuremath{\mathfrak{m}}\textunderscore{}1 to 'donc M\textunderscore{}H→M\textunderscore{}1 est fini'.

Assets P0073-A01--P0073-A03 preserve the proof, corollary, coarse-space proposition, and page ending; P0073-A04 isolates formula (3.9.1).


\par\medskip
\hypertarget{D018-P0074}{}
\pdfbookmark[2]{Authority page 74}{D018-P0074-bookmark}
\noindent{\color{LocatorGray}\footnotesize Authority page 74; collected-paper page 567; original folio 216.}\par\smallskip
Authority page 74; collected-paper volume page 567; original paper folio 216. Printed '-216-' and running locator 'DeRa-74' are excluded from body prose.

The page first equates the scheme M\textunderscore{}n\textasciicircum{}o[1/n] with the stack \ensuremath{\mathfrak{m}}\textunderscore{}n\textasciicircum{}o[1/n] for n≥3, then writes \ensuremath{\mathfrak{m}}\textunderscore{}H\textasciicircum{}o[1/n]=[M\textunderscore{}n\textasciicircum{}o[1/n]/H] before passing to the coarse quotient M\textunderscore{}H\textasciicircum{}o[1/n]=M\textunderscore{}n\textasciicircum{}o[1/n]/H. In the radicial-base-change argument the authority visibly has 'Si p∤n, M\textunderscore{}n⊗F̄\textunderscore{}p, donc ... est lisse', omitting the expected predicate; that source-sic form is frozen and supplied only in English. Section 3.12 distinguishes T̂(E), V̂(E), ℤ̂², and (\ensuremath{\mathbb{A}}\textasciicircum{}f)², and the page ends mid-sentence after the lattice T̂'.

Assets P0074-A01--P0074-A04 preserve the quotient, Hom, Tate-module, kernel, and page-ending statements; P0074-A05 isolates the quotient base-change map.


\par\medskip
\hypertarget{D018-P0075}{}
\pdfbookmark[2]{Authority page 75}{D018-P0075-bookmark}
\noindent{\color{LocatorGray}\footnotesize Authority page 75; collected-paper page 568; original folio 217.}\par\smallskip
Authority page 75; collected-paper volume page 568; original paper folio 217. Printed '-217-' and running locator 'DeRa-75' are excluded from body prose.

The opening completes the equivalence between elliptic curves and curves up to isogeny equipped with a lattice. Section 3.13 interprets a lattice by β⁻¹(ℤ̂²), and β\textunderscore{}n identifies T̂/nT̂ with (ℤ/nℤ)². The left coset K\textbackslash{}Isom(V̂(E\textunderscore{}o),(\ensuremath{\mathbb{A}}\textasciicircum{}f)²) is retained. The source-sic plural 'correspondent' after the singular pair (E\textunderscore{}o,T̂) is frozen in French. Section 3.14 conjugates K to gKg⁻¹ and formula (3.14.1) remains explicitly over Spec(ℚ).

Assets P0075-A01--P0075-A03 preserve the equivalence, level structure, coset, conjugation, and page ending; P0075-A04 isolates formula (3.14.1).


\par\medskip
\hypertarget{D018-P0076}{}
\pdfbookmark[2]{Authority page 76}{D018-P0076-bookmark}
\noindent{\color{LocatorGray}\footnotesize Authority page 76; collected-paper page 569; original folio 218.}\par\smallskip
Authority page 76; collected-paper volume page 569; original paper folio 218. Printed '-218-' and running locator 'DeRa-76' are excluded from body prose.

The page passes to the inverse limit over K and begins the explicit construction of the Hecke correspondence. The inclusions ℤ̂²⊂(1/m)g⁻¹(ℤ̂²)⊂(1/n)ℤ̂² determine B, and the curve is E/α⁻¹(B). The authority's 'Prenons n, et soit m tels que' is frozen in French. Proposition 3.16 extends both g on open stacks and the rational universal-curve isomorphism [g]. The page ends by introducing the finite Hilbert parameter stack I.

Assets P0076-A01--P0076-A04 preserve the inverse-limit action, lattice and subgroup formulas, quotient curve, proposition, and opening proof.


\par\medskip
\hypertarget{D018-P0077}{}
\pdfbookmark[2]{Authority page 77}{D018-P0077-bookmark}
\noindent{\color{LocatorGray}\footnotesize Authority page 77; collected-paper page 570; original folio 219.}\par\smallskip
Authority page 77; collected-paper volume page 570; original paper folio 219. Printed '-219-' and running locator 'DeRa-77' are excluded from body prose.

Lemma 3.17 separates the étale case, controlled by subgroups of E\textunderscore{}b(k), from the non-étale case, where H contains Ker(F) and induction passes to E\textasciicircum{}(p)=E/Ker(F). The proof of 3.16 then prints the fiber product I×\textunderscore{}\{\ensuremath{\mathfrak{m}}\textunderscore{}1\}\ensuremath{\mathfrak{m}}\textunderscore{}K\textasciicircum{}o, without the o-superscript on its base; it extends the section using Zariski's Main Theorem, defines cg by E/B, and obtains an open immersion into the normalization. Applying the construction to g⁻¹ makes it an isomorphism. Section 3.18 states the stack extension, while Proposition 3.19 states the coarse-space extension.

Assets P0077-A01--P0077-A03 preserve the lemma, extension proof, and stack/coarse conclusions; P0077-A04 isolates the final coarse-space isomorphism.


\par\medskip
\hypertarget{D018-P0078}{}
\pdfbookmark[2]{Authority page 78}{D018-P0078-bookmark}
\noindent{\color{LocatorGray}\footnotesize Authority page 78; collected-paper page 571; original folio 220.}\par\smallskip
Authority page 78; collected-paper volume page 571; original paper folio 220. Printed '-220-' and running locator 'DeRa-78' are excluded from body prose.

The graph closure Γ proves Proposition 3.19 using the exact boundary product M\textunderscore{}K\textasciicircum{}∞×\textunderscore{}ℤM\textunderscore{}\{gKg⁻¹\}\textasciicircum{}∞. Section 3.20 prints the quoted name 'e\textunderscore{}n-pairing' and formula e\textunderscore{}n:Λ²C\textunderscore{}n→̃μ\textunderscore{}n. The determinant inverse defines d(α), then the maps to Spec(ℤ[ζ\textunderscore{}n]) and to the det-H invariant subring. The physical page ends on the word 'définissant'; no continuation from page 79 is imported.

Assets P0078-A01--P0078-A03 preserve the graph and cyclotomic construction; P0078-A04 isolates the pairing; P0078-A05 preserves the det-H invariant target and exact page ending.


\par\medskip
\hypertarget{D018-P0079}{}
\pdfbookmark[2]{Authority page 79}{D018-P0079-bookmark}
\noindent{\color{LocatorGray}\footnotesize Authority page 79; collected-paper page 572; original folio 221.}\par\smallskip
Authority page 79; collected-paper volume page 572; original paper folio 221. Printed '-221-' and running locator 'DeRa-79' are excluded from body prose.

Formula (3.20.4) completes the determinant morphism begun at the end of page 78. Section 3.21 works over \ensuremath{\mathfrak{m}}\textunderscore{}\{(m)\}[1/m] with n|m and extends the pairing to the universal generalized curve. The standard Néron m-gon gives formulas (3.21.1)--(3.21.2), with (ℤ/m)\textunderscore{}n identified by i↦(m/n)i; the authority explicitly warns of a possible sign difference. Section 4.1 defines four distinct congruence subgroups, including both Γ\textunderscore{}\{oo\}(n) and Γ'\textunderscore{}\{oo\}(n), and page 79 ends mid-sentence in 4.2.

Assets P0079-A01--P0079-A04 preserve the continuation, pairing, subgroup definitions, and page boundary; P0079-A05 isolates formula (3.21.2).


\par\medskip
\hypertarget{D018-P0080}{}
\pdfbookmark[2]{Authority page 80}{D018-P0080-bookmark}
\noindent{\color{LocatorGray}\footnotesize Authority page 80; collected-paper page 573; original folio 222.}\par\smallskip
Authority page 80; collected-paper volume page 573; original paper folio 222. Printed '-222-' and running locator 'DeRa-80' are excluded from body prose.

The page completes 4.2 and identifies level-Γ\textunderscore{}o(n) data with cyclic subgroup schemes. In 4.4 the authority visibly prints 'le sous-groupe C\textunderscore{}n/A de C/F'; F is undefined in this construction, while every following formula uses the quotient curve C/A. The source-sic C/F is therefore frozen in French and regularized only in English. The involution w sends (C,A) to (C/A,C\textunderscore{}n/A), its square uses C/C\textunderscore{}n→̃C, and its matrix is the one with rows (0,1) and (n,0). Section 4.5 reinterprets w as Cartier duality a↦a\textasciicircum{}*.

Assets P0080-A01--P0080-A03 preserve Constructions 4.3--4.5; P0080-A04 isolates the two involution displays.


\par\medskip
\hypertarget{D018-P0081}{}
\pdfbookmark[2]{Authority page 81}{D018-P0081-bookmark}
\noindent{\color{LocatorGray}\footnotesize Authority page 81; collected-paper page 574; original folio 223.}\par\smallskip
Authority page 81; collected-paper volume page 574; original paper folio 223. Printed '-223-' and running locator 'DeRa-81' are excluded from body prose.

Section 4.6 factors A into its p-primary components and writes w as the composite of the commuting w\textunderscore{}p. Sections 4.7--4.9 distinguish fixing a point from fixing both a point and a quotient class. The exterior-product chain carries C\textunderscore{}n/Im(β) through tensor product and Λ²C\textunderscore{}n to μ\textunderscore{}n; it is the source of the cyclotomic factor in Construction 4.10. Definition 4.11 introduces the leading-prime compactification ′\ensuremath{\mathfrak{m}}\textunderscore{}\{Γ\textunderscore{}o(n)\}, not the unprimed stack.

Assets P0081-A01--P0081-A04 preserve the involutions, exact-order point, cyclotomic factor, and prime-marked definition; P0081-A05 isolates the exterior-product chain.


\par\medskip
\hypertarget{D018-P0082}{}
\pdfbookmark[2]{Authority page 82}{D018-P0082-bookmark}
\noindent{\color{LocatorGray}\footnotesize Authority page 82; collected-paper page 575; original folio 224.}\par\smallskip
Authority page 82; collected-paper volume page 575; original paper folio 224. Printed '-224-' and running locator 'DeRa-82' are excluded from body prose.

The opening parenthesis anticipates Construction 4.13. The forgetful map from ′\ensuremath{\mathfrak{m}}\textunderscore{}\{Γ\textunderscore{}o(n)\}[1/n] has the visibly printed target \ensuremath{\mathfrak{m}}\textunderscore{}* and is étale. The contraction morphism c and automorphism argument establish properness and identify the prime-marked compactification with \ensuremath{\mathfrak{m}}\textunderscore{}\{Γ\textunderscore{}o(n)\}[1/n]. Constructions 4.14--4.15 treat the Γ\textunderscore{}\{oo\} and Γ'\textunderscore{}\{oo\} compactifications. Section 4.16 defines projective level n by the scalar subgroup A and ends on 'Plus généralement :'.

Assets P0082-A01--P0082-A04 preserve the two compactifications, contraction, generalized-level constructions, and exact page ending.


\par\medskip
\hypertarget{D018-P0083}{}
\pdfbookmark[2]{Authority page 83}{D018-P0083-bookmark}
\noindent{\color{LocatorGray}\footnotesize Authority page 83; collected-paper page 576; original folio 225.}\par\smallskip
Authority page 83; collected-paper volume page 576; original paper folio 225. Printed '-225-' and running locator 'DeRa-83' are excluded from body prose.

Definition 4.17 requires every singular geometric fiber to be an n-gon and defines projective level data by Isom(C\textunderscore{}n,(ℤ/n)²)/A. Lemma 4.18 singles out x↦-x. Lemma 4.19 passes from scalar action on X\textunderscore{}n to g²=1: the printed proof uses g-r=n·h, a root of unity ζ, the equality R=ℤ[ζ], and semisimplicity. Construction 4.20 identifies the moduli stack, and Proposition 4.21 begins at the bottom but its map and conclusion are not imported from page 84.

Assets P0083-A01--P0083-A03 preserve the definition, lemmas, and construction; P0083-A04 preserves the exact page-boundary opening of Proposition 4.21.


\par\medskip
\hypertarget{D018-P0084}{}
\pdfbookmark[2]{Authority page 84}{D018-P0084-bookmark}
\noindent{\color{LocatorGray}\footnotesize Authority page 84; collected-paper page 577; original folio 226.}\par\smallskip
Authority page 84; collected-paper volume page 577; original paper folio 226. Printed '-226-' and running locator 'DeRa-84' are excluded from body prose.

The opening map \ensuremath{\mathfrak{m}}\textunderscore{}A[1/n]→M\textunderscore{}A[1/n] and predicate 'est étale' complete Proposition 4.21. Remark 4.22 identifies \ensuremath{\mathfrak{m}}\textunderscore{}A with the leading-prime stack ′\ensuremath{\mathfrak{m}}\textunderscore{}2 at n=2 and retains that prime on ′\ensuremath{\mathfrak{m}}\textunderscore{}2[1/2]→M\textunderscore{}2[1/2]. Section 5 changes to analytic uniformization: the exponential exact sequence carries the injection ι, the period ratio lies in ℂ-ℝ, and the printed matrix convention has rows (a,c) and (b,d). Formula (5.2.1) identifies homology modulo n with E\textunderscore{}n. The page ends after asserting that the class in X\textasciicircum{}±×GL(2,ℤ/n)/GL(2,ℤ) depends only on (E,α).

Assets P0084-A01--P0084-A04 preserve the proposition, analytic sequence, period transformation, and quotient class; P0084-A05 isolates formula (5.2.1).


\par\medskip
\hypertarget{D018-P0085}{}
\pdfbookmark[2]{Authority page 85}{D018-P0085-bookmark}
\noindent{\color{LocatorGray}\footnotesize Authority page 85; collected-paper page 578; original folio 227.}\par\smallskip
Authority page 85; collected-paper volume page 578; original paper folio 227. Printed '-227-' and running locator 'DeRa-85' are excluded from body prose.

Section 5.3 identifies the complex points of the level-H moduli space with analytic double quotients, retaining X in the leftmost quotient, X\textasciicircum{}± in both target quotients, and the finite-adèle group GL(2,\ensuremath{\mathbb{A}}\textasciicircum{}f). Section 5.4 uses surjectivity of SL(2,ℤ)→SL(2,ℤ/n) to identify connected components through determinant map (5.4.1). The cyclotomic map (5.4.2), exponential bijection, and Zariski connectedness yield Corollaries 5.5--5.6.

Assets P0085-A01--P0085-A04 preserve the quotient formulae, determinant/cyclotomic maps, and corollaries; P0085-A05 isolates the quotient-isomorphism displays.


\par\medskip
\hypertarget{D018-P0086}{}
\pdfbookmark[2]{Authority page 86}{D018-P0086-bookmark}
\noindent{\color{LocatorGray}\footnotesize Authority page 86; collected-paper page 579; original folio 228.}\par\smallskip
Authority page 86; collected-paper volume page 579; original paper folio 228. Printed '-228-' and running locator 'DeRa-86' are excluded from body prose.

Section 6 introduces the independently defined prime-marked stack ′\ensuremath{\mathfrak{m}}\textunderscore{}H[1/n], its finite contraction map, and its eventual identification with \ensuremath{\mathfrak{m}}\textunderscore{}H[1/n]. Section 6.1 fixes (A,A',B), β\textunderscore{}0 and β\textunderscore{}1 and realizes them from the generic and special geometric fibers of a generalized elliptic curve over a complete trait. Section 6.2 defines I and Ī and ends with the complete diagram (6.2.1).

Assets P0086-A01--P0086-A04 preserve the section opening, group data, complete-trait example, and exact page ending; P0086-A05 isolates diagram (6.2.1).


\par\medskip
\hypertarget{D018-P0087}{}
\pdfbookmark[2]{Authority page 87}{D018-P0087-bookmark}
\noindent{\color{LocatorGray}\footnotesize Authority page 87; collected-paper page 580; original folio 229.}\par\smallskip
Authority page 87; collected-paper volume page 580; original paper folio 229. Printed '-229-' and running locator 'DeRa-87' are excluded from body prose.

The page defines c:I→Ī and the compatible left GL(2,ℤ/n) and right U actions, then states Lemma 6.2.2. The subgroup correspondence leads to Ī=I/U(B). Adequacy is encoded by the quotient square and the distinct classes a and ā in (6.2.3)--(6.2.4). Lemma 6.2.5 is a disjoint sum over A'⊂B'⊂A. Section 6.3 begins but its isomorphism type continues on page 88.

Assets P0087-A01--P0087-A04 preserve the actions, adequacy conditions, lemma, and exact page boundary; P0087-A05 isolates the quotient square.


\par\medskip
\hypertarget{D018-P0088}{}
\pdfbookmark[2]{Authority page 88}{D018-P0088-bookmark}
\noindent{\color{LocatorGray}\footnotesize Authority page 88; collected-paper page 581; original folio 230.}\par\smallskip
Authority page 88; collected-paper volume page 581; original paper folio 230. Printed '-230-' and running locator 'DeRa-88' are excluded from body prose.

The opening completes the data of section 6.3 and defines \{\}\textunderscore{}H J\textunderscore{}B using left-H and right-Ū freeness. Section 6.4 prescribes the étale sheaf F on smooth fibers and Néron m-gons, with the latter value \{\}\textunderscore{}H J\textunderscore{}B. Section 6.5 begins the concrete systems (B,ᾱ,β) and ends after condition (c); conditions (d)--(e) are confined to page 89.

Assets P0088-A01--P0088-A04 preserve the completed 6.3 data, freeness conditions, sheaf values, and exact 6.5 page boundary.


\par\medskip
\hypertarget{D018-P0089}{}
\pdfbookmark[2]{Authority page 89}{D018-P0089-bookmark}
\noindent{\color{LocatorGray}\footnotesize Authority page 89; collected-paper page 582; original folio 231.}\par\smallskip
Authority page 89; collected-paper volume page 582; original paper folio 231. Printed '-231-' and running locator 'DeRa-89' are excluded from body prose.

Conditions (d)--(e) retain the authority's four labels ᾱ, e\textunderscore{}n, s, and β and its deliberately non-normalized printed top target Λ²ℤ/n in both language layers. In the smooth case the data extend locally to α:C\textunderscore{}n→̃(ℤ/n)\textasciicircum{}2. The two-part equivalence relation yields the presheaf F\textunderscore{}0 and its sheafification F, compatible with base change and represented étale over S. Definition 6.6(i) ends the page; part (ii) is not imported.

Assets P0089-A01--P0089-A04 preserve conditions (d)--(e), the equivalence relation, sheafification, and definition; P0089-A05 isolates the commutative diagram.


\par\medskip
\hypertarget{D018-P0090}{}
\pdfbookmark[2]{Authority page 90}{D018-P0090-bookmark}
\noindent{\color{LocatorGray}\footnotesize Authority page 90; collected-paper page 583; original folio 232.}\par\smallskip
Authority page 90; collected-paper volume page 583; original paper folio 232. Printed '-232-' and running locator 'DeRa-90' are excluded from body prose.

Definition 6.6(ii) defines the prime-marked stack ′\ensuremath{\mathfrak{m}}\textunderscore{}H[1/n]. Forgetting level data is étale; the boundary retains the independent infinity superscript. Theorem 6.7 states properness, smoothness, finite étaleness of the boundary, and representability of c, thereby identifying the prime-marked stack with \ensuremath{\mathfrak{m}}\textunderscore{}H[1/n]. The proof reduces representability to (6.7) and begins the valuative extension construction. The physical page ends after 'on voit que'.

Assets P0090-A01--P0090-A04 preserve the definition, theorem, and valuative construction; P0090-A05 isolates map (6.7).


\par\medskip
\hypertarget{D018-P0091}{}
\pdfbookmark[2]{Authority page 91}{D018-P0091-bookmark}
\noindent{\color{LocatorGray}\footnotesize Authority page 91; collected-paper page 584; original folio 233.}\par\smallskip
Authority page 91; collected-paper volume page 584; original paper folio 233. Printed '-233-' and running locator 'DeRa-91' are excluded from body prose.

This sparse page completes the valuative-extension argument from physical page 90 and proves injectivity of (6.7). The kernel consists of automorphisms acting trivially on Pic⁰(C) and respecting the level structure; condition 6.5(e) forces that kernel to be \{1\}. No Chapter V text is imported.

Assets P0091-A01--P0091-A03 preserve the exact continuation, injectivity proof, and full nonblank page region.


\par\medskip
\hypertarget{D018-P0092}{}
\pdfbookmark[2]{Authority page 92}{D018-P0092-bookmark}
\noindent{\color{LocatorGray}\footnotesize Authority page 92; collected-paper page 585; original folio 234.}\par\smallskip
Authority page 92; collected-paper volume page 585; original paper folio 234. Printed '-234-' and running locator 'DeRa-92' are excluded from body prose.

Chapter V begins the reduction-mod-p study for \ensuremath{\mathfrak{m}}\textunderscore{}K. Section 1 defines the stacks \ensuremath{\mathcal{G}}\textunderscore{}n and ℬ\textunderscore{}n. Proposition 1.2 represents the forgetful map by a formally étale Hilbert scheme. The French source spelling 'reservirons' is retained as printed and separately flagged. The page ends inside the sentence asserting that B acts freely on C.

Assets P0092-A01--P0092-A05 preserve the chapter opening, section introduction, both stack definitions, proposition, proof, and exact boundary.


\par\medskip
\hypertarget{D018-P0093}{}
\pdfbookmark[2]{Authority page 93}{D018-P0093-bookmark}
\noindent{\color{LocatorGray}\footnotesize Authority page 93; collected-paper page 586; original folio 235.}\par\smallskip
Authority page 93; collected-paper volume page 586; original paper folio 235. Printed '-235-' and running locator 'DeRa-93' are excluded from body prose.

The exact sequence 0→B→C\textunderscore{}n→A→0 yields map (1.3.1). Proposition 1.4 identifies ℬ\textunderscore{}n with \ensuremath{\mathcal{G}}\textunderscore{}n. Its proof constructs the inverse from a marked degree-n étale covering and invokes Lemma 1.5. Tildes, smooth-locus superscripts, and the β-labelled isomorphism are distinct marks. The final sentence is left syntactically open.

Assets P0093-A01--P0093-A05 preserve the sequence, proposition, covering construction, lemma, and exact ending.


\par\medskip
\hypertarget{D018-P0094}{}
\pdfbookmark[2]{Authority page 94}{D018-P0094-bookmark}
\noindent{\color{LocatorGray}\footnotesize Authority page 94; collected-paper page 587; original folio 236.}\par\smallskip
Authority page 94; collected-paper volume page 587; original paper folio 236. Printed '-236-' and running locator 'DeRa-94' are excluded from body prose.

The page completes Proposition 1.4, states Theorem 1.6, and introduces the provisional prime-marked stack ′\ensuremath{\mathfrak{m}}\textunderscore{}\{Γ\textunderscore{}0(p)\}. Its identification with the established stack is reduced to assertions A--E: algebraicity, properness, finiteness/representability, normality, and density. Source punctuation and the parenthetical 'groupes(commutatifs)' are preserved.

Assets P0094-A01--P0094-A05 preserve the completed proof, theorem, provisional definition, and the full A--E list.


\par\medskip
\hypertarget{D018-P0095}{}
\pdfbookmark[2]{Authority page 95}{D018-P0095-bookmark}
\noindent{\color{LocatorGray}\footnotesize Authority page 95; collected-paper page 588; original folio 237.}\par\smallskip
Authority page 95; collected-paper volume page 588; original paper folio 237. Printed '-237-' and running locator 'DeRa-95' are excluded from body prose.

The two points at infinity e\textunderscore{}1 and e\textunderscore{}2 come respectively from the 1-gon and p-gon. The open substacks \ensuremath{\mathcal{U}}\textunderscore{}1 and \ensuremath{\mathcal{U}}\textunderscore{}2 separate irreducible fibers from smooth-or-p-gon fibers. Lemma 1.7 supplies representability. Formula (1.8.1) is transcribed with its printed C\textunderscore{}p/A\textunderscore{}p denominator and is not silently normalized.

Assets P0095-A01--P0095-A05 preserve both boundary sections, the open-cover argument, Lemma 1.7, and formula (1.8.1).


\par\medskip
\hypertarget{D018-P0096}{}
\pdfbookmark[2]{Authority page 96}{D018-P0096-bookmark}
\noindent{\color{LocatorGray}\footnotesize Authority page 96; collected-paper page 589; original folio 238.}\par\smallskip
Authority page 96; collected-paper volume page 589; original paper folio 238. Printed '-238-' and running locator 'DeRa-96' are excluded from body prose.

Lemma 1.8.2 is supported by a four-node open-cover diagram. Construction 1.9 glues w and w⁻¹ to an involution of the provisional stack. Proposition 1.10 gives smoothness away from supersingular characteristic-p points. The properness proof begins with the valuative criterion and states Lemma 1.11 in full.

Assets P0096-A01--P0096-A05 preserve the lemma, diagram, involution, proposition, and valuative statement; P0096-A06 isolates the diagram.


\par\medskip
\hypertarget{D018-P0097}{}
\pdfbookmark[2]{Authority page 97}{D018-P0097-bookmark}
\noindent{\color{LocatorGray}\footnotesize Authority page 97; collected-paper page 590; original folio 239.}\par\smallskip
Authority page 97; collected-paper volume page 590; original paper folio 239. Printed '-239-' and running locator 'DeRa-97' are excluded from body prose.

The page completes the existence and uniqueness proof of Lemma 1.11, then begins the proof that c is finite and representable. It treats the boundary fibers and the ordinary finite-distance case, retaining the exact sequence 0→μ\textunderscore{}p→C\textunderscore{}p→ℤ/p→0 and the p+1/2 point count.

Assets P0097-A01--P0097-A05 preserve both valuative paragraphs, the part-(4) transition, the exact sequence, and the terminal characteristic split.


\par\medskip
\hypertarget{D018-P0098}{}
\pdfbookmark[2]{Authority page 98}{D018-P0098-bookmark}
\noindent{\color{LocatorGray}\footnotesize Authority page 98; collected-paper page 591; original folio 240.}\par\smallskip
Authority page 98; collected-paper volume page 591; original paper folio 240. Printed '-240-' and running locator 'DeRa-98' are excluded from body prose.

The infinitesimal α\textunderscore{}\{p²\} case completes the finite-fiber argument. Lemma 1.12 asserts that c is finite flat; Lemma 1.13 supplies the constant-geometric-rank criterion, proved by Nakayama's lemma over a reduced local ring. The visibly empty locus between 'sur' and 'que' is marked ⟦blanc dans l'autorité⟧ in the French layer rather than silently supplied.

Assets P0098-A01--P0098-A05 preserve the α\textunderscore{}\{p²\} argument, automorphism injection, both lemmas, freeness proof, and open page boundary after Spec(k)\ensuremath{\amalg}μ\textunderscore{}p.


\par\medskip
\hypertarget{D018-P0099}{}
\pdfbookmark[2]{Authority page 99}{D018-P0099-bookmark}
\noindent{\color{LocatorGray}\footnotesize Authority page 99; collected-paper page 592; original folio 241.}\par\smallskip
Authority page 99; collected-paper volume page 592; original paper folio 241. Printed '-241-' and running locator 'DeRa-99' are excluded from body prose.

Rank-p subgroup schemes inside α\textunderscore{}\{p²\} are expressed by additive polynomials. The authority's division identity gives a\textasciicircum{}\{p+1\}=0 and T≃Spec(k[a]/(a\textasciicircum{}\{p+1\})). The remainder derives assertion (D) from Cohen--Macaulayness, smoothness off a codimension-two finite set, and Serre's criterion. The unprimed stack at infinity and the printed reference '1 10' are retained exactly.

Assets P0099-A01--P0099-A05 preserve the page continuation, both displayed equations, the full polynomial identity, the local model T, and the normality proof.


\par\medskip
\hypertarget{D018-P0100}{}
\pdfbookmark[2]{Authority page 100}{D018-P0100-bookmark}
\noindent{\color{LocatorGray}\footnotesize Authority page 100; collected-paper page 593; original folio 242.}\par\smallskip
Authority page 100; collected-paper volume page 593; original paper folio 242. Printed '-242-' and running locator 'DeRa-100' are excluded from body prose.

Section 1.14 records variants for full and Γ\textunderscore{}0 level; the second label is printed '(.14.2)' and remains so. Section 1.15 constructs Φ\textunderscore{}1 from Frobenius and Φ\textunderscore{}2=wΦ\textunderscore{}1 from Verschiebung, then gives the c/cw diagram and identifies its composites.

Assets P0100-A01--P0100-A05 preserve both variants, the two objects, the full diagram, and the exact Frobenius-composite ending.


\par\medskip
\hypertarget{D018-P0101}{}
\pdfbookmark[2]{Authority page 101}{D018-P0101-bookmark}
\noindent{\color{LocatorGray}\footnotesize Authority page 101; collected-paper page 594; original folio 243.}\par\smallskip
Authority page 101; collected-paper volume page 594; original paper folio 243. Printed '-243-' and running locator 'DeRa-101' are excluded from body prose.

Theorem 1.16 identifies the two characteristic-p components, their transverse supersingular intersections, the resulting ordinary double points, and regularity. The proof uses the h-open complement and the differential tests dc and d(wc). Variant 1.17 begins and ends immediately before its object descriptions.

Assets P0101-A01--P0101-A05 preserve all theorem clauses, the h-open isomorphism, the transversality proof, and the exact open ending.


\par\medskip
\hypertarget{D018-P0102}{}
\pdfbookmark[2]{Authority page 102}{D018-P0102-bookmark}
\noindent{\color{LocatorGray}\footnotesize Authority page 102; collected-paper page 595; original folio 244.}\par\smallskip
Authority page 102; collected-paper volume page 595; original paper folio 244. Printed '-244-' and running locator 'DeRa-102' are excluded from body prose.

The page supplies the Γ(n) and Γ(n)∩Γ\textunderscore{}0(p) functors promised on page 101, including the symmetric p-isogeny description. It lists Φ\textunderscore{}1, w, and c, warns that w² multiplies the level structure by p, and draws the Frobenius diagram with its vertical pα\textunderscore{}n arrow. The second object list lacks a printed superscript o, and F = E/H) has an unmatched parenthesis. The terminal formula visibly has α\textunderscore{}r; α\textunderscore{}n is the likely contextual correction. The final supersingular condition is syntactically incomplete.

Assets P0102-A01--P0102-A05 preserve both functors, all three arrows, the w² warning, the complete diagram, and the exact terminal fragment.


\par\medskip
\hypertarget{D018-P0103}{}
\pdfbookmark[2]{Authority page 103}{D018-P0103-bookmark}
\noindent{\color{LocatorGray}\footnotesize Authority page 103; collected-paper page 596; original folio 245.}\par\smallskip
Authority page 103; collected-paper volume page 596; original paper folio 245. Printed '-245-' and running locator 'DeRa-103' are excluded from body prose.

The top square completes the level-n supersingular condition begun on page 102. Variant 1.18 gives Frobenius gluing, sections 1.19--1.20 prove and vary regularity, and section 2 opens with the normalized Γ′\textunderscore{}∞(p) moduli problem. The local equation retains uppercase X,Y and the printed multiplication period.

Assets P0103-A01--P0103-A05 preserve the mandatory dense square, Frobenius gluing, the completed-local-ring argument, section transition, and a native 1937x2870 full-page fallback.


\par\medskip
\hypertarget{D018-P0104}{}
\pdfbookmark[2]{Authority page 104}{D018-P0104-bookmark}
\noindent{\color{LocatorGray}\footnotesize Authority page 104; collected-paper page 597; original folio 246.}\par\smallskip
Authority page 104; collected-paper volume page 597; original paper folio 246. Printed '-246-' and running locator 'DeRa-104' are excluded from body prose.

The opening display completes 2.1, after which \ensuremath{\mathfrak{m}}\textunderscore{}\{Γ′\textunderscore{}∞(p)\} is defined by normalization. Sections 2.2--2.4 describe the ordinary opens and their Isom covers, glue the h-open part, and introduce Λ\textunderscore{}p through the character χ. The open-locus glyphs are \ensuremath{\mathcal{G}}′\textunderscore{}p and ℬ′\textunderscore{}p.

Assets P0104-A01--P0104-A05 preserve the normalization formula, the \ensuremath{\mathcal{G}}′\textunderscore{}p/ℬ′\textunderscore{}p and \ensuremath{\mathcal{U}}/\ensuremath{\mathcal{V}} displays, Proposition 2.3, condition (2.4.1), and a native full-page fallback.


\par\medskip
\hypertarget{D018-P0105}{}
\pdfbookmark[2]{Authority page 105}{D018-P0105-bookmark}
\noindent{\color{LocatorGray}\footnotesize Authority page 105; collected-paper page 598; original folio 247.}\par\smallskip
Authority page 105; collected-paper volume page 598; original paper folio 247. Printed '-247-' and running locator 'DeRa-105' are excluded from body prose.

The page completes the description of Spec(Λ\textunderscore{}p), constructs the χ-eigenspace line L of the affine algebra \ensuremath{\mathcal{O}}(A), records Cartier duality and a.b=w(χ), and gives the classification by triples (L,a,b). Printed agreement irregularities and the final L/L* tension are retained and noted.

Assets P0105-A01--P0105-A05 preserve the exact opening continuation, all tensor-power relations, the equivalence of categories, examples (2.4.2)--(2.4.4), and a native full-page fallback.


\par\medskip
\hypertarget{D018-P0106}{}
\pdfbookmark[2]{Authority page 106}{D018-P0106-bookmark}
\noindent{\color{LocatorGray}\footnotesize Authority page 106; collected-paper page 599; original folio 248.}\par\smallskip
Authority page 106; collected-paper volume page 599; original paper folio 248. Printed '-248-' and running locator 'DeRa-106' are excluded from body prose.

Section 2.5 adjoins ζ and derives w(χ,ζ) and the coordinates x,y. Section 2.6 defines the auxiliary left-prime stack, (2.6.1) compares it generically with the normalized stack, and Theorem 2.7 extends that isomorphism. The ordinary-locus proof stops exactly after '(y,ζ)'.

Assets P0106-A01--P0106-A05 preserve the ζ equations, quadruple stack, generic and integral isomorphisms, open proof boundary, and a native full-page fallback.


\par\medskip
\hypertarget{D018-P0107}{}
\pdfbookmark[2]{Authority page 107}{D018-P0107-bookmark}
\noindent{\color{LocatorGray}\footnotesize Authority page 107; collected-paper page 600; original folio 249.}\par\smallskip
Authority page 107; collected-paper volume page 600; original paper folio 249. Printed '-249-' and running locator 'DeRa-107' are excluded from body prose.

The proof of Theorem 2.7 reduces normality to the supersingular points. Lemma 2.8 gives both the strict-henselian local equation and the normalized x,y presentation, after which Remark 2.9 states versality and Variant 2.10 begins a direct modular description.

Assets P0107-A01--P0107-A05 preserve the proof continuation, the mandatory dense Lemma 2.8 box, its local-ring proof, Variant 2.10, and a native full-page fallback.


\par\medskip
\hypertarget{D018-P0108}{}
\pdfbookmark[2]{Authority page 108}{D018-P0108-bookmark}
\noindent{\color{LocatorGray}\footnotesize Authority page 108; collected-paper page 601; original folio 250.}\par\smallskip
Authority page 108; collected-paper volume page 601; original paper folio 250. Printed '-250-' and running locator 'DeRa-108' are excluded from body prose.

The first line completes Variant 2.10. Section 2.11 derives total supersingular ramification, Theorem 2.12 states the two-component and regularity results, and section 2.13 transfers them to full level n by Frobenius gluing. The printed '2.8.montre' is retained, and the final clause remains incomplete after 'les'.

Assets P0108-A01--P0108-A05 preserve the compatibility clause, ramification argument, all of Theorem 2.12, the mandatory dense section 2.13 box, and a native full-page fallback.


\par\medskip
\hypertarget{D018-P0109}{}
\pdfbookmark[2]{Authority page 109}{D018-P0109-bookmark}
\noindent{\color{LocatorGray}\footnotesize Authority page 109; collected-paper page 602; original folio 251.}\par\smallskip
Authority page 109; collected-paper volume page 602; original paper folio 251. Printed '-251-' and running locator 'DeRa-109' are excluded from body prose.

The first line completes section 2.13 from page 108. Sections 2.14--2.15 introduce the intermediate groups Γ\textunderscore{}∞(H), Γ′\textunderscore{}∞(H), the invariant cyclotomic ring, their generic moduli descriptions, and the quotient coordinate u. The printed positive tensor exponent in the last line is retained in French; the coherent negative exponent is disclosed in the translation note.

Assets P0109-A01--P0109-A05 preserve the two matrix groups, the relative-Isom formulas, the quotient-sheaf interpretation, the power-map formula, and a native 1937x2870 full-page fallback.


\par\medskip
\hypertarget{D018-P0110}{}
\pdfbookmark[2]{Authority page 110}{D018-P0110-bookmark}
\noindent{\color{LocatorGray}\footnotesize Authority page 110; collected-paper page 603; original folio 252.}\par\smallskip
Authority page 110; collected-paper volume page 603; original paper folio 252. Printed '-252-' and running locator 'DeRa-110' are excluded from body prose.

Sections 2.16--2.17 define the auxiliary left-prime stack first over Λ\textunderscore{}p and then integrally, with generic quotient data and the u,v equations. Theorem 2.18 gives the comparison, special-fiber components, regularity, and supersingular bijection. The source's prime-less equality in 2.18(i), the printed reference 1.15, and Γ\textunderscore{}0(H) in 2.17 are exposed rather than silently assimilated.

Assets P0110-A01--P0110-A05 preserve the invariant scalar, all tensor-power equations, the leading-prime comparison (2.16.1), the integral data, Theorem 2.18, and a native full-page fallback.


\par\medskip
\hypertarget{D018-P0111}{}
\pdfbookmark[2]{Authority page 111}{D018-P0111-bookmark}
\noindent{\color{LocatorGray}\footnotesize Authority page 111; collected-paper page 604; original folio 253.}\par\smallskip
Authority page 111; collected-paper volume page 604; original paper folio 253. Printed '-253-' and running locator 'DeRa-111' are excluded from body prose.

Section 2.19 gives five consequences for Γ′\textunderscore{}∞(H), including component and supersingular bijections and the quotient \ensuremath{\mathcal{J}}\textunderscore{}n/H. Section 3 begins the good-reduction argument, distinguishes the stack \ensuremath{\mathfrak{m}} from the scheme M over ℚ, defines J(n,H) by Pic\textasciicircum{}o, and states the finite-kernel map (3.1.1).

Assets P0111-A01--P0111-A05 preserve the dense five-part theorem, section transition, both Picard and base-change displays, and a native 1937x2870 full-page fallback.


\par\medskip
\hypertarget{D018-P0112}{}
\pdfbookmark[2]{Authority page 112}{D018-P0112-bookmark}
\noindent{\color{LocatorGray}\footnotesize Authority page 112; collected-paper page 605; original folio 254.}\par\smallskip
Authority page 112; collected-paper volume page 605; original paper folio 254. Printed '-254-' and running locator 'DeRa-112' are excluded from body prose.

Theorem 3.2 treats the cokernel I(n,H). Its proof uses the visibly unprimed Γ(H) and unsubscripted ℤ[ζ], then identifies a toric part of dimension ν-1. Proposition 3.3 applies the Néron-model ranks α\textunderscore{}ab, α\textunderscore{}m, and α\textunderscore{}add. Variant 3.4 switches to field parentheses and unprimed Γ\textunderscore{}∞(H), and its last sentence stops after 'de'.

Assets P0112-A01--P0112-A05 preserve the good-reduction theorem, special-fiber Picard display, the rank calculation, Variant 3.4, the open page ending, and a native full-page fallback.


\par\medskip
\hypertarget{D018-P0113}{}
\pdfbookmark[2]{Authority page 113}{D018-P0113-bookmark}
\noindent{\color{LocatorGray}\footnotesize Authority page 113; collected-paper page 606; original folio 255.}\par\smallskip
Authority page 113; collected-paper volume page 606; original paper folio 255. Printed '-255-' and running locator 'DeRa-113' are excluded from body prose.

The page finishes the finite-kernel construction behind 3.2, then gives Variants 3.5--3.6 by Weil restriction and passage to a K-level quotient. Example 3.7(i) specializes to H=\{±1\}. The prime is retained only in Variant 3.5's two visibly primed Γ′\textunderscore{}∞ formulas; the other infinity-level groups are unprimed.

Assets P0113-A01--P0113-A05 preserve the exact sequence, dense scalar and Weil-restriction formulas, the K-level cokernel, Example 3.7(i), and a native 1937x2870 full-page fallback.


\par\medskip
\hypertarget{D018-P0114}{}
\pdfbookmark[2]{Authority page 114}{D018-P0114-bookmark}
\noindent{\color{LocatorGray}\footnotesize Authority page 114; collected-paper page 607; original folio 256.}\par\smallskip
Authority page 114; collected-paper volume page 607; original paper folio 256. Printed '-256-' and running locator 'DeRa-114' are excluded from body prose.

Examples 3.7(ii)--(iii) give the real-quadratic and square-free-level cases. The proof uses the unramified extension ℚ(ζ\textunderscore{}n)\textasciicircum{}\{H\textunderscore{}p\}/ℚ(ζ\textunderscore{}n)\textasciicircum{}H and compares the B and C cokernels by H-invariants. The terminal clause remains open after 'quotient de'.

Assets P0114-A01--P0114-A05 preserve both examples, the mandatory subgroup matrix, the prime-sum cokernel, the dense B/C comparison, and a native 1937x2870 full-page fallback.


\par\medskip
\hypertarget{D018-P0115}{}
\pdfbookmark[2]{Authority page 115}{D018-P0115-bookmark}
\noindent{\color{LocatorGray}\footnotesize Authority page 115; collected-paper page 608; original folio 257.}\par\smallskip
Authority page 115; collected-paper volume page 608; original paper folio 257. Printed '-257-' and running locator 'DeRa-115' are excluded from body prose.

The opening sentence closes the good-reduction argument from page 114. Part 3.7(iv) permits -1∈H and passage from ℚ(ζ\textunderscore{}p)\textasciicircum{}H to its totally real subfield; Remark 3.8 identifies Casselman's examples. Section 4 then begins the study of \ensuremath{\mathfrak{m}}\textunderscore{}n. Section 4.1 fixes the n|m and p|n fiber conditions, the e\textunderscore{}n exterior pairing, and determinant-one isomorphisms through the displayed commutative square. The page ends exactly at the colon introducing the three cases of 4.2.

Assets P0115-A01--P0115-A05 preserve the opening continuation, the equality of moduli spaces, the 4.1 hypotheses and pairing, the determinant-one square, and a native 1937x2870 full-page fallback.


\par\medskip
\hypertarget{D018-P0116}{}
\pdfbookmark[2]{Authority page 116}{D018-P0116-bookmark}
\noindent{\color{LocatorGray}\footnotesize Authority page 116; collected-paper page 609; original folio 258.}\par\smallskip
Authority page 116; collected-paper volume page 609; original paper folio 258. Printed '-258-' and running locator 'DeRa-116' are excluded from body prose.

The page supplies all three cases announced in 4.2: constant level, μ\textunderscore{}n×ℤ/n level, and the push-out R(A). The printed b.a / b\textasciicircum{}a tension is preserved in French and disclosed in the translation note. The push-out square and exterior-power chain lead into 4.3's determinant-one description of \ensuremath{\mathfrak{m}}\textunderscore{}n[1/n]. Section 4.4 introduces \ensuremath{\mathfrak{a}} and ends at the complete comparison (4.4.1).

Assets P0116-A01--P0116-A05 preserve the three cases, the push-out square, the exterior-power identifications, the moduli descriptions in 4.3--4.4, and a native full-page fallback.


\par\medskip
\hypertarget{D018-P0117}{}
\pdfbookmark[2]{Authority page 117}{D018-P0117-bookmark}
\noindent{\color{LocatorGray}\footnotesize Authority page 117; collected-paper page 610; original folio 259.}\par\smallskip
Authority page 117; collected-paper volume page 610; original paper folio 259. Printed '-259-' and running locator 'DeRa-117' are excluded from body prose.

The page completes the condition defining \ensuremath{\mathcal{G}}′, constructs γ : \ensuremath{\mathcal{G}}′ → ℬ\textunderscore{}n, and proves the smoothness statement of Proposition 4.5 through an explicit determinant-one fiber functor. Section 4.6 then factors n as p\textasciicircum{}k·m and begins the p-primary analysis with (4.6.1). The printed 4.3.3 reference is preserved alongside the visible (4.4.3) label.

Assets P0117-A01--P0117-A05 preserve the defining condition, morphism and proposition, determinant diagram, p-primary transition, and a native 1937x2870 full-page fallback.


\par\medskip
\hypertarget{D018-P0118}{}
\pdfbookmark[2]{Authority page 118}{D018-P0118-bookmark}
\noindent{\color{LocatorGray}\footnotesize Authority page 118; collected-paper page 611; original folio 260.}\par\smallskip
Authority page 118; collected-paper volume page 611; original paper folio 260. Printed '-260-' and running locator 'DeRa-118' are excluded from body prose.

Parts (b)--(f) separate a level-n structure into p\textasciicircum{}k- and m-components, including the cyclotomic compatibility convention and determinant criterion. Lemma 4.7 gives the smoothness of \ensuremath{\mathcal{G}}, and 4.8 begins the projective-line indexing by cyclic subgroups. The printed '(4.6.2) (4.6.4)' cross-reference is retained as an authority tension.

Assets P0118-A01--P0118-A05 preserve both decomposition displays, the compatibility paragraph, factorized moduli descriptions, Lemma 4.7, the exact open ending of 4.8(a), and a native full-page fallback.


\par\medskip
\hypertarget{D018-P0119}{}
\pdfbookmark[2]{Authority page 119}{D018-P0119-bookmark}
\noindent{\color{LocatorGray}\footnotesize Authority page 119; collected-paper page 612; original folio 261.}\par\smallskip
Authority page 119; collected-paper volume page 612; original paper folio 261. Printed '-261-' and running locator 'DeRa-119' are excluded from body prose.

The page completes 4.8(b), identifies the smooth stack C\textunderscore{}A, and records (4.8.1). It then removes the finitely many supersingular points from \ensuremath{\mathfrak{m}}\textunderscore{}1, proves the quasi-finiteness setup, glues the local stacks into C, and opens the valuative proof that c:C→\ensuremath{\mathfrak{m}}\textunderscore{}1\textasciicircum{}h is finite and separated. The terminal clause remains open after 'trait complet à'.

Assets P0119-A01--P0119-A05 preserve both over-tilde arrows, formula (4.8.1), the definitions and lemmas, the gluing formulas, the exact page ending, and a native 1937x2870 full-page fallback.


\par\medskip
\hypertarget{D018-P0120}{}
\pdfbookmark[2]{Authority page 120}{D018-P0120-bookmark}
\noindent{\color{LocatorGray}\footnotesize Authority page 120; collected-paper page 613; original folio 262.}\par\smallskip
Authority page 120; collected-paper volume page 613; original paper folio 262. Printed '-262-' and running locator 'DeRa-120' are excluded from body prose.

The page completes the valuative criterion, extends the generalized elliptic curve and its level structure, and uses the multiplicative p\textasciicircum{}k-subgroup to prove Lemma 4.11.3. It then defines the normalization \ensuremath{\mathfrak{m}}\textunderscore{}n\textasciicircum{}h, identifies it with C in Theorem 4.12, states smoothness in Corollary 4.13, and begins 4.14. The printed A≃μ\textunderscore{}n and B≃ℤ/n are retained despite the surrounding C\textunderscore{}\{p\textasciicircum{}k\} notation.

Assets P0120-A01--P0120-A05 preserve the dense-open and finite-map formulas, the α′′ extension, the exact sequence and A-indexed uniqueness, the theorem and corollary, the exact opening of 4.14, and a native 1937x2870 full-page fallback.


\par\medskip
\hypertarget{D018-P0121}{}
\pdfbookmark[2]{Authority page 121}{D018-P0121-bookmark}
\noindent{\color{LocatorGray}\footnotesize Authority page 121; collected-paper page 614; original folio 263.}\par\smallskip
Authority page 121; collected-paper volume page 614; original paper folio 263. Printed '-263-' and running locator 'DeRa-121' are excluded from body prose.

The page completes the opening of 4.14, resolves a level-p\textasciicircum{}k isomorphism into the exact extension square (4.14.1), and records the equivalent trivialization and component data. Section 4.15 introduces \ensuremath{\mathcal{J}}\textunderscore{}k\textasciicircum{}h and the fiber-product square (4.15.1); Lemma 4.16(i) begins and is deliberately cut after 'la flèche verticale gauche'.

Assets P0121-A01--P0121-A05 preserve the three over-tilde vertical arrows, τ and α′\textunderscore{}1/α′\textunderscore{}2, the two stack squares, the Isom formula, the exact split page boundary, and a native 1937x2870 full-page fallback.


\par\medskip
\hypertarget{D018-P0122}{}
\pdfbookmark[2]{Authority page 122}{D018-P0122-bookmark}
\noindent{\color{LocatorGray}\footnotesize Authority page 122; collected-paper page 615; original folio 264.}\par\smallskip
Authority page 122; collected-paper volume page 615; original paper folio 264. Printed '-264-' and running locator 'DeRa-122' are excluded from body prose.

The page completes Lemma 4.16, states the iterated-Frobenius and torsor descriptions, normalizes \ensuremath{\mathfrak{m}}\textunderscore{}1 in the Igusa covering, and gives the complete-ramification theorem. It then forms \ensuremath{\mathcal{G}}\textasciicircum{}o by fiber product, determines the irreducible components, and states complete ramification of \ensuremath{\mathfrak{m}}\textunderscore{}n→\ensuremath{\mathfrak{m}}\textunderscore{}m. Printed tensions \ensuremath{\mathcal{J}}\textunderscore{}n\textasciicircum{}h/\ensuremath{\mathcal{J}}\textunderscore{}k\textasciicircum{}h, ℬ\textunderscore{}n/ℬ\textunderscore{}\{p\textasciicircum{}k\}, 'corollaire'/Théorème 4.19, the barred/unbarred fields, and the duplicate 4.20 labels are retained.

Assets P0122-A01--P0122-A05 preserve all three parts of Lemma 4.16, Theorem 4.17, formula (4.18.1), Theorem 4.19, both printed 4.20 labels, and a native 1937x2870 full-page fallback.


\par\medskip
\hypertarget{D018-P0123}{}
\pdfbookmark[2]{Authority page 123}{D018-P0123-bookmark}
\noindent{\color{LocatorGray}\footnotesize Authority page 123; collected-paper page 616; original folio 265.}\par\smallskip
Authority page 123; collected-paper volume page 616; original paper folio 265. Printed '-265-' and running locator 'DeRa-123' are excluded from body prose.

The page continues the proof of 4.20 by passing to a supersingular formal deformation Ŝ and the normalization T̂ carrying p\textasciicircum{}k-level structure. Irreducible components are indexed by P, while connected components define the distinct partition \ensuremath{\mathcal{P}}. Lemma 4.21 then begins and compares two supersingular deformation spaces through a prime-to-p isogeny, ending with the induced φ\textunderscore{}T. The printed compact E\textunderscore{}\{oN\} notation is retained without interpretive expansion.

Assets P0123-A01--P0123-A05 preserve the p-primary reduction, formal normalization, component partition, isogeny-deformation diagram, exact final sentence, and a native 1937x2870 full-page fallback.


\par\medskip
\hypertarget{D018-P0124}{}
\pdfbookmark[2]{Authority page 124}{D018-P0124-bookmark}
\noindent{\color{LocatorGray}\footnotesize Authority page 124; collected-paper page 617; original folio 266.}\par\smallskip
Authority page 124; collected-paper volume page 617; original paper folio 266. Printed '-266-' and running locator 'DeRa-124' are excluded from body prose.

The page completes Lemma 4.21 by describing the component bijection through the generic kernel of u and its compatibility with φ\textunderscore{}T. It then proves 4.20: a nontrivial block X of \ensuremath{\mathcal{P}} would make the union of the closed strata (\ensuremath{\mathcal{C}}\textunderscore{}A ⊗ F̄\textunderscore{}p)⁻ a connected component and force \ensuremath{\mathfrak{m}}\textunderscore{}n ⊗ F̄\textunderscore{}p to be disconnected, contrary to (IV 5.5). The lower blank field is not treated as omitted text.

Assets P0124-A01--P0124-A05 preserve the morphism u, component diagram, closure-union contradiction, complete authority body, and a native 1937x2870 full-page fallback.


\par\medskip
\hypertarget{D018-P0125}{}
\pdfbookmark[2]{Authority page 125}{D018-P0125-bookmark}
\noindent{\color{LocatorGray}\footnotesize Authority page 125; collected-paper page 618; original folio 267.}\par\smallskip
Authority page 125; collected-paper volume page 618; original paper folio 267. Printed '-267-' and running locator 'DeRa-125' are excluded from body prose.

This page opens Chapter VI, distinguishes the moduli stack \ensuremath{\mathfrak{m}}\textunderscore{}H from its coarse space M\textunderscore{}H, states Theorem 1.1, and begins its proof through parts a) and b). The authority-printed conjunction 'et' in the theorem is retained rather than emended. The open substack G, its étale image A, the inclusion G⊃\ensuremath{\mathfrak{m}}\textunderscore{}1\textasciicircum{}∞, and the barred geometric fiber (M\textunderscore{}1)\textunderscore{}\{s̄\} remain explicit. The page ends exactly when reducedness is obtained.

Assets P0125-A01--P0125-A05 preserve the chapter transition, theorem and formula (1.1.1), both proof parts, the automorphism condition, the exact physical-page ending, and a native 1937x2870 full-page fallback.


\par\medskip
\hypertarget{D018-P0126}{}
\pdfbookmark[2]{Authority page 126}{D018-P0126-bookmark}
\noindent{\color{LocatorGray}\footnotesize Authority page 126; collected-paper page 619; original folio 268.}\par\smallskip
Authority page 126; collected-paper volume page 619; original paper folio 268. Printed '-268-' and running locator 'DeRa-126' are excluded from body prose.

This page begins with part c), completes the proof of Theorem 1.1, constructs the j-isomorphism from the Néron section and the cohomology of \ensuremath{\mathcal{O}}(M\textunderscore{}1\textasciicircum{}∞), and supplies two classical proofs of (1.1.1). The first uses the Poincaré half-plane; the second uses the invariant-differential parameter scheme, its universal Weierstrass equation, and the weight-(4,6) \ensuremath{\mathbb{G}}\textunderscore{}m quotient. The printed singular 'démonstration' after 'deux' remains visible in the French layer.

Assets P0126-A01--P0126-A05 preserve the continuation boundary, cohomology and j-map, both proofs, the universal equation, the \ensuremath{\mathbb{G}}\textunderscore{}m action, the complete page ending, and a native 1937x2870 full-page fallback.


\par\medskip
\hypertarget{D018-P0127}{}
\pdfbookmark[2]{Authority page 127}{D018-P0127-bookmark}
\noindent{\color{LocatorGray}\footnotesize Authority page 127; collected-paper page 620; original folio 269.}\par\smallskip
Authority page 127; collected-paper volume page 620; original paper folio 269. Printed '-269-' and running locator 'DeRa-127' are excluded from body prose.

The page derives supersingular finiteness from geometric irreducibility, fixes the classical j-map through Tate's formulary, and records the exceptional automorphism sections where the stack-to-coarse-space map is not étale. The missing punctuation in paragraph 1.4 and the local dotted/undotted forms of 2\textasciicircum{}6 3\textasciicircum{}3 remain visible. The last sentence is intentionally incomplete at the physical page boundary.

Assets P0127-A01--P0127-A05 preserve the corollary, Frobenius argument, j-factorization, exceptional automorphisms, Lemma 1.5, ramification, exact page-edge break, and a native 1937x2870 authority fallback usable for both reader layers.


\par\medskip
\hypertarget{D018-P0128}{}
\pdfbookmark[2]{Authority page 128}{D018-P0128-bookmark}
\noindent{\color{LocatorGray}\footnotesize Authority page 128; collected-paper page 621; original folio 270.}\par\smallskip
Authority page 128; collected-paper volume page 621; original paper folio 270. Printed '-270-' and running locator 'DeRa-128' are excluded from body prose.

The page completes the no-section assertion, gives Tate's two families over the punctured j-line, computes their invariants and discriminants, and classifies their twists by an étale ℤ/2-torsor. Proposition 1.7 then supplies elliptic curves at the remaining exceptional invariant values. The reused letter k, the compressed normalization phrase, and the authority's inconsistent multiplication-dot typography are recorded without emendation.

Assets P0128-A01--P0128-A05 preserve the continuation, base S, both Weierstrass equations, invariant calculations, torsor classification, Proposition 1.7, terminal equations, exact page ending, and a native 1937x2870 authority fallback usable for both reader layers.


\par\medskip
\hypertarget{D018-P0129}{}
\pdfbookmark[2]{Authority page 129}{D018-P0129-bookmark}
\noindent{\color{LocatorGray}\footnotesize Authority page 129; collected-paper page 622; original folio 271.}\par\smallskip
Authority page 129; collected-paper volume page 622; original paper folio 271. Printed '-271-' and running locator 'DeRa-129' are excluded from body prose.

The page closes §1 with the automorphism calculation of Lemma 1.9, opens §2 with the criterion for an algebraic stack to be an algebraic space, and begins Lemma 2.2 on lifting polarized abelian varieties with automorphisms. The printed number-theoretic and deformation notation is retained, as are the singular 'variété' and opening 'Soient' tensions. The final display is preserved literally as Spec(W/k)[[t\textunderscore{}1...t\textunderscore{}r]] followed by a dimension formula using k, despite the surrounding prose's W(k). The page ends exactly at that deformation-space formula.

Assets P0129-A01--P0129-A05 are bounded within the native 1937x2870 authority and support raw and presentation review of both language layers.


\par\medskip
\hypertarget{D018-P0130}{}
\pdfbookmark[2]{Authority page 130}{D018-P0130-bookmark}
\noindent{\color{LocatorGray}\footnotesize Authority page 130; collected-paper page 623; original folio 272.}\par\smallskip
Authority page 130; collected-paper volume page 623; original paper folio 272. Printed '-272-' and running locator 'DeRa-130' are excluded from body prose.

The page completes Lemma 2.2 via the fixed formal subscheme M\textasciicircum{}G, proves Proposition 2.3 by a cyclotomic-degree argument, and states Corollaries 2.4 and 2.5. The printed conclusion 'donne le lemme' and Corollary 2.4's self-reference to 2.4 remain explicit. English supplies a minimal continuation cue but does not import omitted source prose.

Assets P0130-A01--P0130-A05 are bounded within the native 1937x2870 authority and support raw and presentation review of both language layers.


\par\medskip
\hypertarget{D018-P0131}{}
\pdfbookmark[2]{Authority page 131}{D018-P0131-bookmark}
\noindent{\color{LocatorGray}\footnotesize Authority page 131; collected-paper page 624; original folio 273.}\par\smallskip
Authority page 131; collected-paper volume page 624; original paper folio 273. Printed '-273-' and running locator 'DeRa-131' are excluded from body prose.

This page gives item 2.6, Theorem 2.7 and its proof, then opens section 3 and paragraph 3.1. The authority's self-citation 'D'après 2.6', the distinction between \ensuremath{\mathfrak{m}}\textunderscore{}H\textasciicircum{}o and M\textunderscore{}H\textasciicircum{}o, barred and unbarred alpha forms, ℚ[ζ\textunderscore{}n]\textasciicircum{}\{det H\}, β, z(β), and the exact physical-page ending are preserved. The printed agreement form 'ᾱ est défini' is recorded without emendation.

Assets P0131-A01--P0131-A05 provide bounded raw crops and presentation derivatives for the source and English layers, including the theorem criterion, determinant field, homology/stabilizer calculation, section transition, and a native 1937x2870 full-page fallback.


\par\medskip
\hypertarget{D018-P0132}{}
\pdfbookmark[2]{Authority page 132}{D018-P0132-bookmark}
\noindent{\color{LocatorGray}\footnotesize Authority page 132; collected-paper page 625; original folio 274.}\par\smallskip
Authority page 132; collected-paper volume page 625; original paper folio 274. Printed '-274-' and running locator 'DeRa-132' are excluded from body prose.

This page completes 3.1, proves Proposition 3.2 through Case 1, and proves Lemma 3.3. The underlined Aut, Giraud H\textasciicircum{}2, obstruction classes, μ\textunderscore{}2/μ\textunderscore{}4/μ\textunderscore{}6 alternatives, statement of (i)--(ii), and every term and label in the long exact sequence remain explicit. Two printed tensions are not repaired: (ii) has target Gal(k̄/k,μ\textunderscore{}n) without H\textasciicircum{}2, and the final quotient display repeats k\textasciicircum{}*/k\textasciicircum{}\{*m/n\} on both sides.

Assets P0132-A01--P0132-A05 provide bounded raw crops and presentation derivatives for the source and English layers, including the descent obstruction, Proposition 3.2, Case 1, Lemma 3.3, the exact sequences and quotient map, and a native 1937x2870 full-page fallback.


\par\medskip
\hypertarget{D018-P0133}{}
\pdfbookmark[2]{Authority page 133}{D018-P0133-bookmark}
\noindent{\color{LocatorGray}\footnotesize Authority page 133; collected-paper page 626; original folio 275.}\par\smallskip
Authority page 133; collected-paper volume page 626; original paper folio 275. Printed '-275-' and running locator 'DeRa-133' are excluded from body prose.

The page closes the positive-characteristic descent argument through a finite-field reduction, a Giraud H\textasciicircum{}2 obstruction, and a Frobenius-power calculation. It then begins the numerical remarks and defines the Euler--Poincaré characteristic of an algebraic stack by a weighted coarse-moduli stratification. The marked arrows, dotted exceptional value, printed 'résoud', φ′ after scalar extension, unmatched parenthesis in 'Soit \ensuremath{\mathcal{Z}}\textunderscore{}i)', and unindexed defining sum remain visible rather than being silently regularized.

Assets P0133-A01--P0133-A05 preserve the finite-field argument, descent equations, Section 4 transition, definition 4.1, exact page ending, and a native 1937x2870 authority fallback usable for both reader layers.


\par\medskip
\hypertarget{D018-P0134}{}
\pdfbookmark[2]{Authority page 134}{D018-P0134-bookmark}
\noindent{\color{LocatorGray}\footnotesize Authority page 134; collected-paper page 627; original folio 276.}\par\smallskip
Authority page 134; collected-paper volume page 627; original paper folio 276. Printed '-276-' and running locator 'DeRa-134' are excluded from body prose.

The page states additivity, multiplicativity, algebraic extension, and the finite-flat degree formula for the Euler--Poincaré characteristic. It then gives the concrete groupoid formula and begins the application to \ensuremath{\mathfrak{m}}\textunderscore{}1\textasciicircum{}o over ℂ. The source's multiplication period, prose/display Ã\textunderscore{}x versus A\textunderscore{}x tension, missing cardinality bars in (4.1.3), letter-o superscript, and terminal 'auto-' remain exposed.

Assets P0134-A01--P0134-A05 preserve all four properties, displays (4.1.1)--(4.1.3), the opening of paragraph 4.2, exact mid-word page edge, and a native 1937x2870 authority fallback usable for both reader layers.


\par\medskip
\hypertarget{D018-P0135}{}
\pdfbookmark[2]{Authority page 135}{D018-P0135-bookmark}
\noindent{\color{LocatorGray}\footnotesize Authority page 135; collected-paper page 628; original folio 277.}\par\smallskip
Authority page 135; collected-paper volume page 628; original paper folio 277. Printed '-277-' and running locator 'DeRa-135' are excluded from body prose.

The page begins within a sentence from page 134, completes the Euler-characteristic calculation for the open elliptic-moduli stack, passes to the coarse modular curve by accounting for cusps and automorphisms, and starts the definition of degree on a one-dimensional proper Cohen-Macaulay algebraic stack. The visibly printed addition of -1/2, ζ(-1), ±U, product factors, φ(n), and tilded local ring are preserved. The page ends exactly in §4.3(b), after the negative local-order formula.

Assets P0135-A01--P0135-A05 are bounded within the native 1937x2870 authority and support raw and presentation review of both language layers.


\par\medskip
\hypertarget{D018-P0136}{}
\pdfbookmark[2]{Authority page 136}{D018-P0136-bookmark}
\noindent{\color{LocatorGray}\footnotesize Authority page 136; collected-paper page 629; original folio 278.}\par\smallskip
Authority page 136; collected-paper volume page 629; original paper folio 278. Printed '-278-' and running locator 'DeRa-136' are excluded from body prose.

The page completes the degree definition, records its additivity, covering, and specialization properties, computes deg(ω)=1/24 from the discriminant, and begins a direct deformation-theoretic proof of (4.4.2). The asterisk in f*ω is retained exactly as a visible character because its typographic elevation is not determinable from the monospaced scan. All other dense formula symbols are visually resolved. The page ends mid-sentence at 'se prolonge'; no matter from page 137 is imported.

Assets P0136-A01--P0136-A05 are bounded within the native 1937x2870 authority and support raw and presentation review of both language layers.


\par\medskip
\hypertarget{D018-P0137}{}
\pdfbookmark[2]{Authority page 137}{D018-P0137-bookmark}
\noindent{\color{LocatorGray}\footnotesize Authority page 137; collected-paper page 630; original folio 279.}\par\smallskip
Authority page 137; collected-paper volume page 630; original paper folio 279. Printed '-279-' and running locator 'DeRa-137' are excluded from body prose.

This page completes the preceding discussion at infinity, gives (4.5.1)--(4.5.2), derives the Euler-characteristic and degree identities, announces the Hasse-invariant argument, and begins section 4.6. Fraktur \ensuremath{\mathfrak{m}} versus roman M, superscript letter o, every tensor exponent, the Frobenius triangle, and the incomplete Cartier-trace sentence are preserved. The visible overstrike on the running word 'Hasse' is disclosed without treating the word as deleted.

Assets P0137-A01--P0137-A05 provide bounded raw crops and presentation derivatives for both reader layers, including all displays, the section transition, the Frobenius diagram, and a native 1937x2870 full-page fallback.


\par\medskip
\hypertarget{D018-P0138}{}
\pdfbookmark[2]{Authority page 138}{D018-P0138-bookmark}
\noindent{\color{LocatorGray}\footnotesize Authority page 138; collected-paper page 631; original folio 280.}\par\smallskip
Authority page 138; collected-paper volume page 631; original paper folio 280. Printed '-280-' and running locator 'DeRa-138' are excluded from body prose.

This page completes the Cartier trace construction, gives (4.6.1)--(4.6.3), records the dual description of the Hasse invariant and its simple supersingular zeros, and begins the deformation-theoretic paragraph 4.7. The simultaneous f superscript/subscript placement is linearized as f\textunderscore{}*\textasciicircum{}(p) where printed; the later identification sentence's visible omission of the star in R\textasciicircum{}1 f\textasciicircum{}(p)\ensuremath{\mathcal{O}} is retained in French and disclosed as emended in English. The source's A\textunderscore{}o is retained as a letter o, and the unfinished final sentence is not completed from another page.

Assets P0138-A01--P0138-A05 provide bounded raw crops and presentation derivatives for both reader layers, including the trace and pullback maps, finite-field assertions, de Rham definition, exact sequence, and a native 1937x2870 full-page fallback.


\par\medskip
\hypertarget{D018-P0139}{}
\pdfbookmark[2]{Authority page 139}{D018-P0139-bookmark}
\noindent{\color{LocatorGray}\footnotesize Authority page 139; collected-paper page 632; original folio 281.}\par\smallskip
Authority crops P0139-A01--A04 cover the continuation formula, both diagrams, the deformation argument, and the supersingular weighted sum; P0139-A05 preserves the full page and both transition edges.

The folio header -281- and locator DeRa-139 are copy matter and are excluded from the transcription. The opening formula and closing clause are retained as incomplete continuations exactly where the authority page cuts them.

Potentially fragile readings were resolved from pixels: A\textunderscore{}o and E\textunderscore{}o use the letter o; the first square has f\textunderscore{}o ⊗ k[ε] above and f* below; the second diagram has F* ⊗ k[ε] in the middle vertical arrow and a capital Σ in the kernel argument.


\par\medskip
\hypertarget{D018-P0140}{}
\pdfbookmark[2]{Authority page 140}{D018-P0140-bookmark}
\noindent{\color{LocatorGray}\footnotesize Authority page 140; collected-paper page 633; original folio 282.}\par\smallskip
Authority crops P0140-A01--A04 cover formula (4.9.1), the examples, the Section 5 transition, the ± U matrix, and both quotient descriptions; P0140-A05 preserves the full page and both transition edges.

The folio header -282- and locator DeRa-140 are copy matter and are excluded from the transcription. The page begins with the equality introduced on page 139 and ends at the colon introducing a scheme isomorphism on page 141.

Potentially fragile readings were resolved from pixels: the general count has (1-p)/2; the heading says pointes; the automorphism image is ± U; and the Section 5.2 display is a left quotient by H and a right quotient by ± U.


\par\medskip
\hypertarget{D018-P0141}{}
\pdfbookmark[2]{Authority page 141}{D018-P0141-bookmark}
\noindent{\color{LocatorGray}\footnotesize Authority page 141; collected-paper page 634; original folio 283.}\par\smallskip
Authority page 141; collected-paper volume page 634; original paper folio 283. Printed '-283-' and running locator 'DeRa-141' are excluded from body prose.

The page gives Construction 5.3, begins the underlined section 6, passes from the stack of cyclic-kernel isogenies to its coarse modular curve, defines the maps c and cw, and constructs the symmetric modular equation Φ\textunderscore{}n. Fraktur \ensuremath{\mathfrak{m}} versus roman M, every lowercase letter o, the double quotient by H and ±U, the coefficient range a,b≤P(n), and both forms n≠1 and n>1 are preserved. The printed underlines beneath Construction 5.3, the section heading, and 'l'équation modulaire' remain visible in the authority assets rather than becoming editorial markup.

Assets P0141-A01--P0141-A05 are bounded within the native 1937x2870 authority and cover every display, section transition, coefficient assertion, and the exact continuation into page 142.


\par\medskip
\hypertarget{D018-P0142}{}
\pdfbookmark[2]{Authority page 142}{D018-P0142-bookmark}
\noindent{\color{LocatorGray}\footnotesize Authority page 142; collected-paper page 635; original folio 284.}\par\smallskip
Authority page 142; collected-paper volume page 635; original paper folio 284. Printed '-284-' and running locator 'DeRa-142' are excluded from body prose.

The page opens with the normalized symmetric expansion (6.3.2), explains the possible noninjectivity of the modular map through two cyclic-kernel isogenies, derives complex multiplication from g=e\textasciicircum{}\{-1\}φ\textunderscore{}2\textasciicircum{}*fφ\textunderscore{}1, gives the converse quotient construction, and reaches Proposition 6.5. The strict coefficient bound, star placement, degree n\textasciicircum{}2, Ker(g)≠E\textunderscore{}n, opposite quotient arrows, p∤n, stack/coarse typography, and the underlined proposition statement are retained. The visibly different e and f arrows are disclosed rather than silently harmonized.

Assets P0142-A01--P0142-A05 are bounded within the native 1937x2870 authority and cover all formulas, arrow displays, Proposition 6.5, and the unfinished final sentence without importing page-143 matter.


\par\medskip
\hypertarget{D018-P0143}{}
\pdfbookmark[2]{Authority page 143}{D018-P0143-bookmark}
\noindent{\color{LocatorGray}\footnotesize Authority page 143; collected-paper page 636; original folio 285.}\par\smallskip
Authority page 143; collected-paper volume page 636; original paper folio 285. Printed '-285-' and running locator 'DeRa-143' are excluded from body prose.

This page completes the preceding n-isogeny sentence, states Theorem 6.6, records congruences (6.6.1)--(6.6.4), proves the smoothness setup of Proposition 6.7 away from the exceptional sections, and opens the characteristic-2 local argument. Letter-o Γ\textunderscore{}o(p), primes, p-powers, stack versus coarse glyphs, the printed underlining, and the unfinished lower boundary are preserved or explicitly disclosed.

Assets P0143-A01--P0143-A05 provide bounded raw crops and presentation derivatives for all formulas and transitions, plus a native 1937x2870 full-page authority fallback.


\par\medskip
\hypertarget{D018-P0144}{}
\pdfbookmark[2]{Authority page 144}{D018-P0144-bookmark}
\noindent{\color{LocatorGray}\footnotesize Authority page 144; collected-paper page 637; original folio 286.}\par\smallskip
Authority page 144; collected-paper volume page 637; original paper folio 286. Printed '-286-' and running locator 'DeRa-144' are excluded from body prose.

This page completes the characteristic-2 strict-henselization argument, gives the characteristic-3 analogue, begins section 6.8, and states all three parts of Theorem 6.9. Every K∩Γ\textunderscore{}o(p) subscript uses letter o. The special-fiber gluing and the final two-case formula for k are retained exactly; printed underlines are treated as emphasis, not correction marks.

Assets P0144-A01--P0144-A05 cover the continuation, finite-field and genus-zero arguments, section 6.8, Theorem 6.9, and the exact final piecewise display, with a native 1937x2870 full-page fallback.


\par\medskip
\hypertarget{D018-P0145}{}
\pdfbookmark[2]{Authority page 145}{D018-P0145-bookmark}
\noindent{\color{LocatorGray}\footnotesize Authority page 145; collected-paper page 638; original folio 287.}\par\smallskip
Authority crops P0145-A01--A04 cover the local equation, proof, coarse-moduli diagram, and Corollary 6.10; P0145-A05 preserves the full page. Folio -287- and locator DeRa-145 are excluded as copy matter. The diagram's arrow topology and the stack/coarse-space distinction were independently verified from the authority pixels.


\par\medskip
\hypertarget{D018-P0146}{}
\pdfbookmark[2]{Authority page 146}{D018-P0146-bookmark}
\noindent{\color{LocatorGray}\footnotesize Authority page 146; collected-paper page 639; original folio 288.}\par\smallskip
Authority crops P0146-A01--A04 cover the definition of A, both quotient models, the branch restriction, and the divisor argument; P0146-A05 preserves the full page. Folio -288- and locator DeRa-146 are excluded as copy matter. The authority visibly prints [[X'Y']] without a comma; this likely error is preserved, not silently repaired.


\par\medskip
\hypertarget{D018-P0147}{}
\pdfbookmark[2]{Authority page 147}{D018-P0147-bookmark}
\noindent{\color{LocatorGray}\footnotesize Authority page 147; collected-paper page 640; original folio 289.}\par\smallskip
Authority crops P0147-A01--A04 cover both Euler formulas, the resolution reminder, Lemma 6.14, and Corollary 6.15; P0147-A05 preserves the full page. Folio -289- and locator DeRa-147 are excluded as copy matter. The page break after classes d'isomorphie de is semantic and preserved.


\par\medskip
\hypertarget{D018-P0148}{}
\pdfbookmark[2]{Authority page 148}{D018-P0148-bookmark}
\noindent{\color{LocatorGray}\footnotesize Authority page 148; collected-paper page 641; original folio 290.}\par\smallskip
Authority crops P0148-A01--A04 cover the continued corollary, supersingularity criteria, the complete incidence diagram, and denominator conclusion; P0148-A05 preserves the full page. Folio -290- and locator DeRa-148 are excluded as copy matter. P0148-A03 is the required image fallback for exact component incidence.


\par\medskip
\hypertarget{D018-P0149}{}
\pdfbookmark[2]{Authority page 149}{D018-P0149-bookmark}
\noindent{\color{LocatorGray}\footnotesize Authority page 149; collected-paper page 642; original folio 291.}\par\smallskip
Authority crops P0149-A01--A04 cover the chapter opening and equations (1.1.1)--(1.1.6); P0149-A05 preserves the full page. Folio -291- and locator DeRa-149 are excluded as copy matter. The isolated oblique pixel in (1.1.1) is treated as a scan blemish, and the source's S[1/t]. ; punctuation is retained.


\par\medskip
\hypertarget{D018-P0150}{}
\pdfbookmark[2]{Authority page 150}{D018-P0150-bookmark}
\noindent{\color{LocatorGray}\footnotesize Authority page 150; collected-paper page 643; original folio 292.}\par\smallskip
Authority crops P0150-A01--A04 cover the continuation formulas, examples, infinite-chain schematic, and Néron-model conclusion; P0150-A05 preserves the full page. Folio -292- and locator DeRa-150 are excluded as copy matter. The printed y\textunderscore{}i=t\textasciicircum{}\{-i\}y\textunderscore{}0 anomaly is retained, and P0150-A03 controls the chain and brace topology.


\par\medskip
\hypertarget{D018-P0151}{}
\pdfbookmark[2]{Authority page 151}{D018-P0151-bookmark}
\noindent{\color{LocatorGray}\footnotesize Authority page 151; collected-paper page 644; original folio 293.}\par\smallskip
Authority crops P0151-A01--A04 cover the section comparison, group law, action, and involution; P0151-A05 preserves the full page. Folio -293- and locator DeRa-151 are excluded as copy matter. The visible superscript T\textasciicircum{}k is anomalous beside T\textunderscore{}i but is preserved, and the iso-/morphe split is exact.


\par\medskip
\hypertarget{D018-P0152}{}
\pdfbookmark[2]{Authority page 152}{D018-P0152-bookmark}
\noindent{\color{LocatorGray}\footnotesize Authority page 152; collected-paper page 645; original folio 294.}\par\smallskip
Authority crops P0152-A01--A04 cover [u], the contraction, quotient construction, and Example 1.7; P0152-A05 preserves the full page. Folio -294- and locator DeRa-152 are excluded as copy matter. Barred and unbarred multiplicative-group models remain distinct, and the final définit continues to page 153.


\par\medskip
\hypertarget{D018-P0153}{}
\pdfbookmark[2]{Authority page 153}{D018-P0153-bookmark}
\noindent{\color{LocatorGray}\footnotesize Authority page 153; collected-paper page 646; original folio 295.}\par\smallskip
Authority crops P0153-A01--A04 cover the opening display, Proposition 1.9, formal limit, and algebraization; P0153-A05 preserves the full page. Folio -295- and locator DeRa-153 are excluded as copy matter. Direct pixel inspection resolves the recurring glyph as barred Ḡ\textunderscore{}m\textasciicircum{}t; the literal dot and right-pointing limit arrow are retained.


\par\medskip
\hypertarget{D018-P0154}{}
\pdfbookmark[2]{Authority page 154}{D018-P0154-bookmark}
\noindent{\color{LocatorGray}\footnotesize Authority page 154; collected-paper page 647; original folio 296.}\par\smallskip
Authority crops P0154-A01--A04 cover nonsmoothness, formal groups, Lie and torsion maps, and (1.13.1)--(1.13.2); P0154-A05 preserves the full page. Folio -296- and locator DeRa-154 are excluded as copy matter. Hats, bars, injection arrows, source agreement errors, and the terminal semicolon were independently checked.


\par\medskip
\hypertarget{D018-P0155}{}
\pdfbookmark[2]{Authority page 155}{D018-P0155-bookmark}
\noindent{\color{LocatorGray}\footnotesize Authority page 155; collected-paper page 648; original folio 297.}\par\smallskip
Authority crops P0155-A01--A04 cover (1.13.3)--(1.13.4), both contractions, and Construction 1.15; P0155-A05 preserves the full page. Folio -297- and locator DeRa-155 are excluded as copy matter. The radical carries an upper-left k index, and the source's missing closing parenthesis after Spec(A[t\textasciicircum{}\{-1\}] is retained.


\par\medskip
\hypertarget{D018-P0156}{}
\pdfbookmark[2]{Authority page 156}{D018-P0156-bookmark}
\noindent{\color{LocatorGray}\footnotesize Authority page 156; collected-paper page 649; original folio 298.}\par\smallskip
Authority crops P0156-A01--A04 cover Definition 1.16, structures (1.16.1)--(1.16.4), the pairing, and Theorem 2.1; P0156-A05 preserves the full page. Folio -298- and locator DeRa-156 are excluded as copy matter. The opening \ensuremath{\mathbb{G}}\textunderscore{}m\textasciicircum{}t is unbarred; the e\textunderscore{}k/e\textunderscore{}n mismatch, extra period, and final continuation comma are preserved.


\par\medskip
\hypertarget{D018-P0157}{}
\pdfbookmark[2]{Authority page 157}{D018-P0157-bookmark}
\noindent{\color{LocatorGray}\footnotesize Authority page 157; collected-paper page 650; original folio 299.}\par\smallskip
Authority crops P0157-A01--A04 cover the Theorem 2.1 continuation, Corollary 2.2, section 2.3, and (2.3.1); P0157-A05 preserves the full page. Folio -299- and locator DeRa-157 are copy matter and excluded. The τ, f, q, ζ\textunderscore{}n, and f\textunderscore{}n boundary are retained.


\par\medskip
\hypertarget{D018-P0158}{}
\pdfbookmark[2]{Authority page 158}{D018-P0158-bookmark}
\noindent{\color{LocatorGray}\footnotesize Authority page 158; collected-paper page 651; original folio 300.}\par\smallskip
Authority crops P0158-A01--A04 cover the q\textasciicircum{}\{1/n\} Tate curve, (2.3.2), commutative square, and completion; P0158-A05 preserves the full page. Folio -300- and DeRa-158 are excluded. The upper-triangular stabilizer ±U is diplomatic source notation.


\par\medskip
\hypertarget{D018-P0159}{}
\pdfbookmark[2]{Authority page 159}{D018-P0159-bookmark}
\noindent{\color{LocatorGray}\footnotesize Authority page 159; collected-paper page 652; original folio 301.}\par\smallskip
Authority crops P0159-A01--A04 cover Corollaries 2.5--2.6, X\textunderscore{}v, \ensuremath{\mathcal{G}}\textunderscore{}v, and §2.7; P0159-A05 preserves the full page. Folio -301- and DeRa-159 are excluded. The local units u and v and the terminal continuation are not harmonized.


\par\medskip
\hypertarget{D018-P0160}{}
\pdfbookmark[2]{Authority page 160}{D018-P0160-bookmark}
\noindent{\color{LocatorGray}\footnotesize Authority page 160; collected-paper page 653; original folio 302.}\par\smallskip
Authority crops P0160-A01--A04 cover the signed j(E) formula, §3, and Lemma 3.3; P0160-A05 preserves the full page. Folio -302- and DeRa-160 are excluded. The printed ± and finite-covering page edge are preserved.


\par\medskip
\hypertarget{D018-P0161}{}
\pdfbookmark[2]{Authority page 161}{D018-P0161-bookmark}
\noindent{\color{LocatorGray}\footnotesize Authority page 161; collected-paper page 654; original folio 303.}\par\smallskip
Authority crops P0161-A01--A04 cover Theorem 3.4, Δ, coefficient definitions, and Tate evaluation; P0161-A05 preserves the full page. Folio -303- and DeRa-161 are excluded. Tensor powers and the Fourier continuation are checked against pixels.


\par\medskip
\hypertarget{D018-P0162}{}
\pdfbookmark[2]{Authority page 162}{D018-P0162-bookmark}
\noindent{\color{LocatorGray}\footnotesize Authority page 162; collected-paper page 655; original folio 304.}\par\smallskip
Authority crops P0162-A01--A04 cover (3.7.1), gφ, (3.8.1), and Theorem 3.9; P0162-A05 preserves the full page. Folio -304- and DeRa-162 are excluded. The finite-type ring and ζ\textunderscore{}n\textasciicircum{}\{amn\} exponent are retained.


\par\medskip
\hypertarget{D018-P0163}{}
\pdfbookmark[2]{Authority page 163}{D018-P0163-bookmark}
\noindent{\color{LocatorGray}\footnotesize Authority page 163; collected-paper page 656; original folio 305.}\par\smallskip
Authority crops P0163-A01–A04 cover the two proofs, exact diagram, and Theorem 3.10; P0163-A05 preserves the full page. Folio -305- and DeRa-163 are copy matter and excluded. The formal-completion notation D̂, coefficient rings, and fₙ boundary are retained.


\par\medskip
\hypertarget{D018-P0164}{}
\pdfbookmark[2]{Authority page 164}{D018-P0164-bookmark}
\noindent{\color{LocatorGray}\footnotesize Authority page 164; collected-paper page 657; original folio 306.}\par\smallskip
Authority crops P0164-A01–A04 cover Corollaries 3.11–3.14 and ν(π,g,φ); P0164-A05 preserves the full page. Folio -306- and DeRa-164 are excluded. The final ‘on a’ routes to the page-165 inequality.


\par\medskip
\hypertarget{D018-P0165}{}
\pdfbookmark[2]{Authority page 165}{D018-P0165-bookmark}
\noindent{\color{LocatorGray}\footnotesize Authority page 165; collected-paper page 658; original folio 307.}\par\smallskip
Authority crops P0165-A01–A04 cover the BSL inequality, Corollary 3.15, §3.16, and the first Tate expansion; P0165-A05 preserves the full page. Folio -307- and DeRa-165 are excluded. The second expansion comma is a physical boundary.


\par\medskip
\hypertarget{D018-P0166}{}
\pdfbookmark[2]{Authority page 166}{D018-P0166-bookmark}
\noindent{\color{LocatorGray}\footnotesize Authority page 166; collected-paper page 659; original folio 308.}\par\smallskip
Authority crops P0166-A01–A04 cover the second expansion, involution w, valuation identities, and divisor setup; P0166-A05 preserves the full page. Folio -308- and DeRa-166 are excluded. Subscripts 1 and 2 and the p-power substitutions are retained.


\par\medskip
\hypertarget{D018-P0167}{}
\pdfbookmark[2]{Authority page 167}{D018-P0167-bookmark}
\noindent{\color{LocatorGray}\footnotesize Authority page 167; collected-paper page 660; original folio 309.}\par\smallskip
Authority crops P0167-A01–A04 cover component intersections and Proposition 3.20; P0167-A05 preserves the full page. Folio -309- and DeRa-167 are excluded. The final ‘et donc’ is not silently completed.


\par\medskip
\hypertarget{D018-P0168}{}
\pdfbookmark[2]{Authority page 168}{D018-P0168-bookmark}
\noindent{\color{LocatorGray}\footnotesize Authority page 168; collected-paper page 661; original folio 310.}\par\smallskip
Authority crops P0168-A01–A04 cover the p-integrality consequence and §§4.1–4.4; P0168-A05 preserves the full page. Folio -310- and DeRa-168 are excluded. Complex uniformization and the final analytic comparison are retained.


\par\medskip
\hypertarget{D018-P0169}{}
\pdfbookmark[2]{Authority page 169}{D018-P0169-bookmark}
\noindent{\color{LocatorGray}\footnotesize Authority page 169; collected-paper page 662; original folio 311.}\par\smallskip
Authority crops P0169-A01-P0169-A04 cover the §4.4 continuation, transformation law, slash operator, cusp condition, and Fourier series; P0169-A05 preserves the full page. Folio -311- and DeRa-169 are excluded. The page starts mid-§4.4 and ends at a complete display.


\par\medskip
\hypertarget{D018-P0170}{}
\pdfbookmark[2]{Authority page 170}{D018-P0170-bookmark}
\noindent{\color{LocatorGray}\footnotesize Authority page 170; collected-paper page 663; original folio 312.}\par\smallskip
Authority crops P0170-A01-P0170-A04 cover Construction 4.6, its proof, §4.7, and the opening of §4.8; P0170-A05 preserves the full page. Folio -312- and DeRa-170 are excluded. The first §4.8 bullet continues without completion on page 171.


\par\medskip
\hypertarget{D018-P0171}{}
\pdfbookmark[2]{Authority page 171}{D018-P0171-bookmark}
\noindent{\color{LocatorGray}\footnotesize Authority page 171; collected-paper page 664; original folio 313.}\par\smallskip
Authority crops P0171-A01-P0171-A04 preserve both terminal §4.8 bullets and their sparse topology; P0171-A05 preserves the full page. Folio -313- and DeRa-171 are excluded. The large blank area is not silently removed or populated.


\par\medskip
\hypertarget{D018-P0172}{}
\pdfbookmark[2]{Authority page 172}{D018-P0172-bookmark}
\noindent{\color{LocatorGray}\footnotesize Authority page 172; collected-paper page 665; original folio 314.}\par\smallskip
Authority crops P0172-A01-P0172-A04 cover bibliography [1]-[15]; P0172-A05 preserves the full page. Folio -314- and DeRa-172 are excluded. Bibliographic entries are source content, while only running copy matter is removed. The printed capital P and typo 'prpperties' in [13] are retained as authority text.


\par\medskip
\hypertarget{D018-P0173}{}
\pdfbookmark[2]{Authority page 173}{D018-P0173-bookmark}
\noindent{\color{LocatorGray}\footnotesize Authority page 173; collected-paper page 666; original folio 315.}\par\smallskip
Authority crops P0173-A01-P0173-A04 cover bibliography [16]-[34]; P0173-A05 preserves the full page. Folio -315- and DeRa-173 are excluded. The authority's absent labels [21] and [28] remain absent.


\par\medskip
\hypertarget{D018-P0174}{}
\pdfbookmark[2]{Authority page 174}{D018-P0174-bookmark}
\noindent{\color{LocatorGray}\footnotesize Authority page 174; collected-paper page 667; original folio 316.}\par\smallskip
Authority crops P0174-A01-P0174-A04 cover references [35]-[36], SIGLES, and the Katz-Mazur note; P0174-A05 preserves the full page. Folio -316- and DeRa-174 are excluded. This is the exact authority terminus.


\par\medskip
\end{document}
